% ============================================================
%  F. C. Sibbern, Om Philosophiens Begreb, Natur og Væsen.
%    En Fremstilling af Philosophiens Propædeutik.
%  Kjøbenhavn: Paa Forfatterens eget Forlag (Trykt hos Directeur
%    Jens Hostrup Schultz), 1843.
%  TRANSCRIPTION (Danish)
%  Source: Royal Danish Library (KB) scan,
%    ~/bibliotek/Sibbern, Frederik/om-philosophiens-begreb.pdf
%  101 PDF pages. Front matter (Fraktur) + Fraktur body (§§ 1--21).
%  Physical extent: [6], 86 s.  Companion piece to Om Erkjendelse og
%    Granskning (1821/22): the Fortale explicitly relates the two.
%  PAGE MAP (verified against printed headers + Indhold):
%    Body (arabic): printed page = PDF page - 11  (§ 1 opens p.1 = PDF 12).
%    Title page = PDF 6; Fortale = PDF 8--9; Indhold + Rettelser = PDF 10--11.
%  RETTELSER (errata, PDF 11): several are SUBSTANTIVE (marked NB) and are to
%    be applied to the Danish, e.g. p.39 Theologie->Teleologie; p.9 "ikke i
%    selv"->"ikke i sig selv"; p.33 Begrebets->Begrebet; p.60 add comma. The
%    full list is transcribed at the end; NB corrections are folded into the
%    body as they arise (marked with a % Rettelse note at each site).
%  See RESUME-NOTES.md for structure, section->page map, and progress.
% ============================================================
\documentclass[12pt,a4paper]{book}
\usepackage[utf8]{inputenc}
\usepackage[T1]{fontenc}
\usepackage{libertinus}
\usepackage{libertinust1math}
\usepackage{textalpha}   % Greek words typed directly where they occur
\usepackage[danish]{babel}
\usepackage{enumitem}
\usepackage[protrusion=true,expansion=true]{microtype}
\usepackage{setspace}
\usepackage[margin=1.2in]{geometry}
\usepackage{fancyhdr}
\setlength{\headheight}{28pt}
\addtolength{\topmargin}{-13.5pt}
\usepackage{hyperref}

\hypersetup{
  pdftitle={Om Philosophiens Begreb, Natur og Væsen},
  pdfauthor={F. C. Sibbern},
  hidelinks
}

\pagestyle{fancy}
\fancyhf{}
\fancyhead[LE,RO]{\thepage}
\fancyhead[CE]{\textit{\leftmark}}
\fancyhead[CO]{\textit{\rightmark}}

\setstretch{1.3}

% Group-heading helper (the centered bold rubrics of the Indhold / body)
\newcommand{\grouphead}[1]{%
  \bigskip\begin{center}\large\textbf{#1}\end{center}\smallskip}

\title{\textbf{Om}\\[6pt]
  \Huge\textbf{Philosophiens Begreb,}\\[6pt]
  \LARGE\textbf{Natur og Væsen.}\\[14pt]
  \normalsize En Fremstilling af\\[4pt]
  \Large\textbf{Philosophiens Propædeutik.}\\[16pt]
  \normalsize Af\\[6pt]
  \large Dr.\ Frederik Christian Sibbern,\\[4pt]
  \normalsize Professor i Philosophien.}
\author{}
\date{Kjøbenhavn, 1843.\\
  Paa Forfatterens eget Forlag\\
  trykt hos Directeur Jens Hostrup Schultz,\\
  Kongelig og Universitets-Bogtrykker.}

\begin{document}
\maketitle
\thispagestyle{empty}
\cleardoublepage

%% ============================================================
%%  FORTALE  (PDF 8--9; Fraktur). Transcribed and image-verified.
%% ============================================================
\chapter*{Fortale.}
\markboth{Fortale}{Fortale}

\begin{center}\rule{0.28\linewidth}{0.4pt}\end{center}

Dersom jeg ikke hidtil havde ladet ethvert Skrivt af mig, der er traadt
selvstændigt frem i Literaturen, ledsages af en Fortale, --- ikke just fordi
jeg udtrykkeligen har henregnet en saadan med til en Bogs rette Udstyrelse,
men fordi det hver Gang er faldet mig naturligt, --- skulde jeg denne Gang
været tilbøielig til at lade al Fortale falde bort, da jeg Intet havde at
sige, som kunde være Gjenstand for en saadan. Jeg vilde nemlig helst undladet
udtrykkelig at bemærke, at, hvad maaskee flere af mine Læsere her ville vente
angivet, vil jeg overlade til dem selv at afgjøre. Jeg mener herved: i hvilket
Forhold nærværende Skrivt staaer til min i Slutningen af Aaret 1821 udgivne
Bog: Om Erkjendelse og Granskning. Skal jeg herover yttre mig, da vidste jeg
med Hensyn til Publicum --- thi hvorledes jeg vil benytte begge Skrivter ved
mine Forelæsninger over Philosophiens Propædeutik, hører det ikke herhen at
forklare --- ikke at sige Andet, end at jeg ønskede, at man i dens Forhold til
nærværende Bog vilde betragte hiin tidligere, som om den var af en anden
Forfatter, hvem jeg havde benyttet, og til hvem jeg derfor oftere havde
henviist, medens jeg dog havde havt en anden Gjenstand for mig, end han. Hiint
er nemlig Videnskab og videnskabeligt Studium overhovedet Gjenstanden, og en
Indledning til det akademiske Studium overhovedet er det, som hiint leveres.
Her derimod har jeg havt Philosophien i Særdeleshed for mig, og har derfor i
Henseende til \emph{denne} havt Mere at gaae ind paa, end hvad der hiint er
fremkommet. --- Tilsidst beder jeg, at blandt Rettelserne, især de, som ere
betegnede med \textbf{NB}, maatte vorde bemærkede.

\bigskip
\noindent\textbf{Kjøbenhavn,} den 28de October 1843.
\nopagebreak
\begin{flushright}\textbf{Sibbern.}\end{flushright}

\begin{center}\rule{0.18\linewidth}{0.4pt}\end{center}

\mainmatter

%% ============================================================
%%  FORELØBIGE BEMÆRKNINGER  (§§ 1--2, printed pp. 1--6; PDF 12--17)
%%  Transcribed and image-verified.
%% ============================================================
\grouphead{Foreløbige Bemærkninger.}

\begin{center}\large\S~1.\end{center}
\markboth{\S~1}{\S~1}
\begin{center}Hvad der her skal foretages, og hvorfra gaaes ud.\end{center}

Det er Philosophiens Begreb, Natur og Væsen, vi her ville søge at opklare os.
Vi kunne da ikke indskrænke os til det Spørgsmaal: Hvad forstaaes ved
Philosophie?, men have med det langt mere sigende Spørgsmaal at gjøre: Hvad er
Philosophie? Imidlertid maae vi dog naturligviis strax og fra først af tage
Hensyn til Ordets Betydning, idet vi agte paa, hvad der kaldes Philosophie og
gjennem Historiens Løb er bleven kaldet saaledes, og under dette Navn har
gjort sig gjeldende. Thi, om end det Begreb, vi skulle søge at give, er at
udvikle og bestemme, som den selve Sagen betænkende Overveielse vil fordre
det, saa kan dog det Spørgsmaal, om nu det ifølge heraf Opstillede svarer til
Det, man kalder Philosophie, eller maaskee, med al den Gyldighed, det i og for
sig monne have, er noget ganske Andet, end hvad man ved Ordet Philosophie har
for Øie, kun afgjøres ved at tage den Betydning, hvormed dette Navn gaaer og
gjelder, i Betragtning. Ved Navnet peges hen paa den Gjenstand, hvorom der
skal tales, og man maa fra først af sørge for, at man ikke kommer til at tale
om en anden Gjenstand, istedetfor om denne. Fra en given Forestilling om, hvad
Philosophie er, en given Betydning af Ordet, gaae vi da rettest ud.

Men selve den Sag, paa hvilken der saaledes ved Navnet er peget hen, og som
derved er hævet ud for os, maae vi nu derpaa vistnok betragte ganske, som den
selv vil sees, og maae forfølge den i dens hele Conseqvents, saa at altsaa
Begrebet længere hen ikke vil kunne bindes til, og i sin hele Gjennemførelse
betinges af, hiin givne Forestilling.

Herved komme vi da nu vistnok til at betragte Philosophien i Philosophiens
eget Lys, og komme altsaa til at lægge vor Philosophie til Grund; thi vi
skulle nu philosophere os til den Erkjendelse, vi søge; men i al Philosophering
gjør Philosophie sig strax gjeldende. Men dette Cykliske er en Følge af Sagens
hele Stilling; thi Philosophien er kun et Led med i den Verdenstotalitet, som
netop Philosophien, ifølge det Begreb om samme, vi maae forudsætte som givet,
skal opklare for os.

{\small\emph{Anm.~1.} Der opstaaer her naturligen det Spørgsmaal, om ikke
Philosophiens Begreb skulde være at afhandle paa to ganske forskjellige
Steder: først i de til Philosophien indledende Betragtninger, hvor det
philosophiske Foretagende jo ikke vil kunne omhandles, uden at et Begreb om,
hvad Philosophie er, lægges til Grund, --- dernæst midt inde i Philosophiens
System, naar den speculative Grunderkjendelse saavidt er tilvunden og uddannet,
at Philosophien, idet den træder frem som et medindgribende Led i den i
Bevidstheden sig udviklende aandelige Leven, bliver at betragte fra det
speculative --- nærmest et psycho-biologiskt --- Standpunct. Man seer let, at,
skjønt Philosophien vil føre os tilbage til alle Tings første Grundmomenter og
Begyndelse (\textit{initia}), saa hører den dog selv hjemme midt inde i
Verdensløbet, hvor den først kan opstaae i den høiere Bevidsthedsleven, eller
midt i Aandelighedens Sphære, og da maa vise sig som et væsentligt Hovedmoment
i den aandelige Levens Totalitet. Spørgsmaalet om Philosophiens egentlige Væsen
hører da hjemme der, hvor ogsaa Poesie, Moralitet, Religiøsitet, skulle
constituere sig for os i deres væsentlige Medhenhøren til Totaliteten af den
aandelige Leven. ---\par}

{\small Der spørges da naturligviis, om vi ikke fra denne aandsphilosophiske og
egentlig speculative Overveielse af Philosophiens Væsen skulle adskille en
anden foreløbig, til en blot Indledning til Philosophien henhørende, hvilken vi
da kunde kalde den blot propædeutiske, der da tillige var at kalde den blot
explicative. Denne skulde da blot forfølge det givne Begreb om Philosophie i al
dets Conseqvents og derefter bestemme og exponere det, medens hiin derimod
skulde lade Philosophien selv constituere sig for os og derved vise sig i sit
egentlige Væsen. Denne sidste vilde da være den, som, naar den meget
udtryksfulde Benævnelse Philosophiens Philosophie skal optages, maatte belægges
med dette Navn.\par}

{\small Men til en saadan blot propædeutisk Exposition af Philosophiens Begreb
kunne og ville vi dog her ikke indskrænke os, eller holde den sondret fra den
speculative, da det er Philosophien selv, --- ikke det Begreb, der har dannet
sig om den, --- vi her have for os, og dennes Opklaring, naar den skal være
fuldelig, ei kan blive staaende ved hiint blotte Propædeutiske.\par}

{\small Imidlertid maae vi dog vistnok bestandigen søge at blive os bevidste,
hvorvidt vort Begreb, alt som det mere og mere bestemmer og udvikler sig for
os, forbliver i det blot Propædeutiskes Gebeet, saa at det kan opstilles som
gyldigt med Hensyn til de forskjellige philosophiske Partier (Platonikere,
Skeptikere, Kantianere, Lockianere o.\,s.\,v.), og altsaa i Philosophiens
Historie kan lægges til Grund; hvorvidt det derimod gaaer videre, og altsaa
allerede uddannes ifølge en bestemt philosophisk Grundanskuelses Tilsagn.\par}

{\small\emph{Anm.~2.} Man vil let see, at det vilde være aldeles Utilsvarende
at give den første og fundamentelle Deel af selve Philosophien, hvilken Hegel
har kaldet Logik, Navn af Philosophiens Philosophie, saaledes som Nogle have
gjort det. Hegel er ogsaa saare langt fra, at have kunnet give nogen Anledning
hertil, da hos ham Philosophiens Philosophie --- dog uden at dette Navn hos ham
findes --- fremkommer i Slutningen af hans hele Systems tredie og sidste
Deel.\par}

{\small\emph{Anm.~3.} Angaaende det Udtryk: \emph{Natur og Væsen}, og dets
Forhold til Udtrykket: \emph{Begreb}, bemærkes, at Det, der her nærmest og
først søges, er Philosophiens Begreb og sammes nøiagtige Bestemmelse, men at
dette Begreb dog dernæst skal udfolde sig for os eller træde ud i det hele
Indbegreb af sin Conseqvents, eller i sit hele Indhold, hvilket da er Det, vi
ved Udtrykket: \emph{Natur og Væsen} have for Øie. (Jeg henviser herved til min
Logik, \S~35.)\par}

\begin{center}\large\S~2.\end{center}
\markboth{\S~2}{\S~2}
\begin{center}Udsigt over det Følgendes hele Gang.\end{center}

Spørge vi da allerførst: Hvad forstaaer man ved Philosophie, hvorved vi kun
foreløbigen skulle angive saa Meget, at vi med Bestemthed kunne vide, hvorom
det er, der skal handles, men ogsaa strax maae komme til at angive de første
Puncter, hvorpaa vi have at henvende vor Opmærksomhed, saa bemærke vi, at det
overleverede Begreb medfører, at Philosophie skal gaae ud paa en vis
Grunderkjendelse, der skal tjene til en dybere Opklaring af hele Tilværelsen,
men at den skal danne sig ved Tænkning og Overveielse, og altsaa blive en
Erkjendelse, som skal være sig sin indre Grundethed bevidst og kunne gjøre Rede
for denne. Men saaledes komme vi da til at lade vor Overveielse gaae følgende
Gang:

Først og fremmest tage vi i Betragtning, at 1) Philosophie er et Indbegreb af
Erkjenden, og overveie derfor Erkjendelsens Natur og Væsen i Almindelighed.

Dernæst bemærke vi, at, eftersom Philosophie gaaer ud paa en vis dybere, til de
sidste eller inderste Grundforhold henskuende Grunderkjendelse, men en saadan
dog kun da viser sig som Philosophie, naar den træder frem som udsprungen af
Fornufterkjendelse, og især kun da, naar Efterforskning eller Randsagelse
derved har været i Bevægelse, saa ledes vi herved til at betragte
Erkjendelsesbestræbelsen i Almindelighed i dens Overgang fra dens første
Hovedtrin, eller den første umiddelbare Erkjendens Trin, til dens andet
Hovedtrin, eller den ved Randsagelsesaanden fremkaldte middelbare, men tillige
høiere, Erkjendens Trin, paa hvilket Erkjendelsen bliver til en Erkjenden
ifølge Indsigt, og tillige --- eftersom herved Prøvelse maa finde Sted,
Regnskab maa aflægges, Grunde komme til Overveielse og Grundethed oplyses og
godtgjøres, --- til en Erkjendelse i Kraft af Grunde, hvilket da lader os see
Philosophien som den der 2) er et Indbegreb af intelligent og rational
Erkjenden.

Men hertil knytter sig da fremdeles den Bemærkning, at da dette Indbegreb maa
danne en velsammenhængende Heelhed, maa Philosophie i sin Gjennemførelse og
Udførelse blive til 3) en videnskabelig Erkjendelse, hvorved vi da bringes til
at betragte Videnskab og Videnskabelighed i Almindelighed.

Men da nu dog --- ifølge det overleverede Begreb --- ikke al videnskabelig
Erkjenden er Philosophie, saa komme vi, ved at spørge os, hvad det nu er for en
Videnskab, hvori Philosophie bestaaer, til at blive vaer, at Bestræbelsen efter
at trænge hen til Grundene ei kan Andet end føre til, at der danner sig, 4) en
\emph{Videnskab}, der paa en vis \emph{altomfattende Maade} gaaer ud paa sidste
Grunde og en sidste Grunderkjendelse, og at det er denne især, man har for Øie,
naar man taler om Philosophie; hvorhos det da, naar nu det Spørgsmaal gjøres,
om det er al Erkjendelses eller den hele Tilværelses sidste Grunde og
Grundforhold, som herved menes, vil vise sig, at, skjønt en dobbelt
Hovedopgave for Philosophien i denne Henseende har dannet sig, og har maattet
danne sig, blive det dog tilsidst kun 5) \emph{Tilværelsens}, eller som man
pleier sige, \emph{alle Tings} sidste Grunde og Grundforhold, ved hvis
Udfindelse og Erkjendelse den altomfattende Grunderkjendelse skal danne sig,
hvori Philosophien bestaaer.

Denne fra gammel Tid af overleverede Bestemmelse af Begrebet Philosophie,
hvilken vi saaledes bringes til at gaae ud fra, viser os Philosophien som den,
der træder frem ligeoverfor og i en vis Modsætning til Indbegrebet af al den
Erkjenden, hvorved alt Det gives os og fremstiller sig for os, om hvis sidste
Grunde og Grundforhold der nu skal være Tale, hvilket vi da saaledes
foranlediges og opfordres til at samle for os i et Overblik.

I sin Forfølgelse i de forskjellige Hovedhenseender eller Hovedretninger ville
vi dernæst see Philosophiens Begreb føre til en heel Cyklus af
Begrebsbestemmelser, der kan uddannes til en heel Cyklus af Definitioner.

Derpaa ville vi da gaae til at betragte Philosophien i sin Udførelse. Men naar
vi da nu allerførst tage i Overveielse, hvad sidste Grunde ville sige, og
hvorledes det forholder sig med vor Hentrængen til dem, vil Spørgsmaalet om
Philosophiens Grundlag og Udgang vise sig at maatte gaae foran for Spørgsmaalet
om Udførelsen af Philosophiens Værk, hvortil da, naar vi saa derpaa
fuldstændigen have besvaret os det Spørgsmaal, hvad Philosophien er, en
Betragtning over Philosophiens Betydning i Totaliteten af Tilværelsen naturligen
vil kunne træde.

Hermed er da den Gang, vi ville tage, heelt angiven.

{\small\emph{Anm.} Det maa her strax bemærkes, hvad der længere hen paa sit
Sted vil blive nærmere omhandlet, at der ved Siden af den egentlige
Philosophie, hvorom her allene er Tale, gives et Indbegreb af Videnskaber, som,
uden at være egentlig Philosophie, dog ogsaa kaldes philosophiske, men som jeg,
til Adskillelse fra hiin, (see længere hen \S~11) giver Navn af explicativ
Philosophie, og her aldeles ikke tager Hensyn til.\par}

\begin{center}\rule{0.18\linewidth}{0.4pt}\end{center}

%% ============================================================
%%  INDLEDENDE BETRAGTNINGER. OM ERKJENDELSE I ALMINDELIGHED.
%%  (§§ 3--7, printed pp. 7--24; PDF 18--35)
%% ============================================================
\grouphead{Indledende Betragtninger.\\[2pt]\normalsize Om Erkjendelse i Almindelighed.}

\begin{center}\large\S~3.\end{center}
\markboth{\S~3}{\S~3}
\begin{center}Philosophie er a) Erkjendelse. --- Al Erkjendelses Grundmomenter.\end{center}

Al Erkjenden gaaer ud paa et \emph{er}, hvorved Noget, idet det gjør sig
gjeldende hos os, statueres af os som et Objectivt, hvilket vi maae tillægge
Gyldighed, og hvilket vi, forsaavidt det træder frem for os i denne sin
Gyldighed, kalde det Sande, der da er at tage ganske saaledes, som det selv
giver sig os, det er: viser sig at være.

Her see vi da en Modsætning imellem en sig selv bevidst Subjectivitet og en
uafhængigt heraf bestaaende Objectivitet, der forresten træder frem for os ei
blot udenfra, men ogsaa indenfra. Men Erkjendelse indtræder først, idet denne
Modsætning saaledes overvindes, at Det, som er, netop ogsaa for Subjectiviteten
træder frem ganske saaledes, som det er, saa at det kommer til for os at være,
det er af os statueres at være, ganske som det er i sig selv, hvorved vi da
komme til, i det Samme at vi leve i vor egen individuelle Subsistents, tillige
at leve i Objectivitetens, tilmed det Sandes, Fylde.

% NB: print reads ``i Eet og og Alt'' --- apparent dittography (``og'' repeated
% at the line break, printed p.8); NOT in the errata leaf; kept as printed.
Det Medierende og Forenende mellem det Sande og Subjectiviteten, der tillige
viser sig som en Individualitet, er fra det Førstes Side den bestemmende
Nødvendighed, hvormed det gjør sig gjeldende, og paa Subjectivitetens Side
\emph{deels} dennes Bestemtvorden og dermed indtrædende Kommen til Bevidsthed,
der da kan kaldes en Modtagen, (\textit{scimus, qvia accepimus}), \emph{deels}
den individuelle Levens egen til hiin Bestemtvorden svarende Gjøren. --- Men
denne sidste er atter deels en saadan, der blot gaaer ud paa en
Sindshenvendelse, deels en saadan, som gaaer ud paa indre subjectiv
Gjenfremstilling og Objectivering. Ved den første see vi en ved den bestemmende
Indflydelses vækkende Magt fremkaldt subjectiv Reaction og Imødetræden, der blot
gaaer ud paa at gjøre Modtagelsen dybere og fuldeligere, saavel ifølge en dermed
forbunden Hengivelse, som ifølge en tiltrædende Besindelse, som Omsorg for, at
det Sande maa i alle sine Puncter og Forhold komme til sin rette fulde Magt hos
os, saa at vi i Eet og og Alt aldeles kunne vorde bestemmede netop i
Overeensstemmelse med det Sandes egen Medfør. I den subjective Gjenfremstilling
og Objectivering see vi derimod en af den individuelle Leven selv udspringende
Productivitet, der har sin indre Virksomhedsgrund og Virksomhedsmaade i sig
selv, og nu, som en frembringende Gjøren, i den subjective Selvbevidsthed
frembringer og constituerer den indre Fremstilling, i og ved hvilket det Sande
kommer til at indoptages i den subjective Leven. (Og her hedder det da egentlig
først: \textit{scimus, qvia facimus}. --- Iøvrigt henviser jeg her til mit
Skrivt „om Erkjend.\ og Granskn." \S\S~1--3, saavelsom til min Psychologie). Men
Det, som gjør, at vi dog i vor egen Productivitet have, ikke noget blot
Subjectivt, men, formedelst og gjennem dette blot Subjective, have det Sande
selv, er den Magt, hvormed dette saaledes bestemmer og behersker den
Forestillingen skabende og dannende subjective Leven, at det herved tillige
fremstiller sig som det Gyldige for den sig selv bevidste Skuen.

{\small\emph{Anm.~1.} Naar der her tales om et „er", saa see vi her ikke allene
en Modsætning imellem et „er" og et „ikke-er", men og en Modsætning imellem et
„er" og et „viser sig" eller „synes". Herved kan der da strax spørges, om da
ikke ogsaa dette „synes" maa siges at være, og omvendt, om et „er" nogensinde
kan være Andet, end et „viser sig". Her vilde da være at spørge, hvad hiint
„er" egentlig vil sige; og naar vi da bemærke, at vi overhovedet kun kunne have
at gjøre med en Væren, forsaavidt den gjør sig gjeldende, men nu dog ogsaa
ethvert Synes gjør sig gjeldende, saa spørges, om her er en forskjellig
Gjørensiggjeldende, og hvorpaa da saadan Forskjellighed beroer; hvorved da vel
en Forskjel imellem en blot momentan og subjectiv og en mere constant og
objectiv Gjørensiggjeldende turde være at antage, og i denne sidste atter være
at adskille a) en blot sporadisk, b) en blot mikrokosmisk, c) en makrokosmisk
Gyldighed, og kun den sidste at være at erkjende for den egentlige sande i
sidste Instants. Men her vilde da nu dog den Betragtning opstaae, at al Skin-,
al sporadisk, al mikrokosmisk Gjørensiggjeldende dog tilsidst maa være grundet i
den makrokosmiske Tingenes Orden, saa at her altsaa var at bemærke, at al hiin
øvrige Gjørensiggjeldende dog kun saaes i sit sande Lys, altsaa i sin egentlige
Væren, ved at sees i sin Fremgaaen af og Tilbageførelighed til den makrokosmiske
Tingenes Orden. --- Men alle de Betragtninger, som saaledes kunde fremkaldes ved
Overveielsen af hiint „er", høre ei herhen, men til selve den speculative
Philosophies Grundbetragtninger. --- Saa Meget er vist, at under al vor
% Rettelse (errata leaf, p.9 l.26): print reads ``gjør sig i Modsætningen'';
% the erratum ``bortfalder Ordet i'' deletes the word ``i''. Applied.
Erkjendelsesstræben gjør sig Modsætningen imellem et „er" og et „synes" --- et
% Rettelse (errata leaf, p.9 l.28): print reads „ikke i selv"; the erratum
% ``for "ikke i selv" læs: "ikke i sig selv"'' inserts ``sig''. Applied.
„i sig selv" og et „ikke i sig selv" --- og tillige med samme Modsætningen
imellem et Objectivt og et Subjectivt, saaledes gjeldende, at vi i vor Erkjendens
hele Gang, saavelsom i vor hele øvrige hermed forbundne Leven, ei kunne unddrage
os den. Og hertil holde vi os her, hvor vi blot have at exponere Erkjendelsens
Natur, som den giver sig os, og hvor vi uden videre tage Existentsen af en
Erkjenden med alle deri liggende Momenter som given, ligesom vi tage Existentsen
af en Philosophie som given, uden at vi her gaae ind i Spørgsmaalet om al
Erkjendelses og al Philosophies egentlige Realitet, hvilken det hører selve den
speculative Philosophie til at afgjøre. --- Det Begreb om Erkjendelse vi her have
givet, har aabenbar --- Dette er endnu at bemærke --- Gyldighed ogsaa for Den,
som vil nægte al vor Erkjenden en sand Realitet; thi ogsaa en Saadan, maa dog
have et Begreb om Erkjendelse, som han lægger til Grund, og det kan kun være det
her givne; og Den, som kun vil tillægge al vor Erkjenden en vis subjectiv
Gyldighed, maa dog tilstaae, at Erkjendelsen beroer paa vor Forestillens
Samstemmen med Noget, som antager Charakteren af et Objectivt.\par}

{\small\emph{Anm.~2.} Idet der danner sig et „er" griber strax Intelligentsen
ind, og en saakaldet Kategorie, eller bedre en apriorisk Grundbestemmelse, gjør
sig gjeldende. Hermed træder da ogsaa strax Erkjendelsen frem som den, der træder
ud i en Dømmen, hvorved da et Noget (som Subject) træder frem som Det, der er
Noget vist (Hvadhedens Grundbestemmelse), og dette Hvad det er, træder da frem
som et Andet, hvilket dog ved selve Dømningen gjøres til Eet med hiint (som
Prædicat). Denne Dømningens Grundnatur udtales i det saakaldte
Identitetsprincip.\par}

{\small\emph{Anm.~3.} Det her opstaaende Spørgsmaal om, hvorledes hiin
Productivitet, naar den har sin indre Virksomhedsgrund og Virksomhedsmaade,
altsaa og Regulativet for denne, i sig selv, dog kommer til at fremstille et
virkeligt Objectivt, er af speculativ Natur, og finder for os sit Svar i Ideen
om en \textit{harmonia originaria} imellem Individualiteten og Totaliteten. Det
Objectives bestemmende Indflydelse er vistnok kun en fremkaldende, da det er den
individuelle Leven selv, som ved sin egen Activitet i Individualitetens
Bevidsthedskreds fremvirker Det, hvorved det Objective skal vorde det Erkjendte.
Men ligesom den hele Individualitet overhovedet først fremkommer ved den totale
Tingenes Gang, og i sin hele Tilvær vorder, som denne fører det med sig,
saaledes constitueres ogsaa dens hele Organisation baade i indvortes sjælelig og
i udvortes Henseende paa en saadan til Totaliteten svarende Maade, at enhver fra
Totaliteten kommende Indvirkning, enten den saa kommer udenfra, eller fra det
indvortes Universelle, fremkalder sin bestemte tilsvarende Activitet i den
individuelle Organisation, eller i Mikrokosmen, og det saa, at den nu, ifølge
dennes egen Grundnatur, netop fører til en saadan vis bestemt Production, at just
hiint bestemte Objective derved forestilles af os. Saa at, uagtet vi bestandig
skue Gjenstandene i vor egen Sjæl, saa skue vi dog de virkelige Gjenstande, som
de i sig selv ere, fordi vore indvortes Forestillinger blive Afbilleder af dem
eller Gjenbilleder til dem, eftersom vort indvortes Bevidsthedsliv, ifølge hiin
sin indre Harmonie med Totaliteten, naar det (vel at mærke) er i sin normale
Gang, hver Gang maa ud af sin egen indre Birken frembringe netop den Nerveaction
med tilsvarende Forestillingsaction, som bevirker netop den Forestilling, der
svarer til Gjenstanden. Den individuelle Organismus, der i sig indeholder den
reale Muelighed til alle muelige Forestillingsproductioner, er ligesom beregnet
til, paa ethvert Punct, d.\ e.\ i enhver Henseende, at virke i en saadan Accord
eller Concordants med enhver fra det Totales Side kommende Indvirkning, at hiin
hver Gang just frembringer, hvad denne gaaer ud paa at fremkalde (som naar et
Instrument ifølge sin egen hele Construction netop, om end paa sin individuelle
Maade, gjengiver just den Tone, der slaaes an af et andet). Vi føle os imidlertid
herunder ikke som i Forhold til blotte Gjenbilleder, men som i Forhold til
virkelige Gjenstande, fordi vi herunder bestandig føle vor Organismus underkastet
en Paavirkning, som, kommende fra noget Andet, end os selv, er den, der fremkalder
og fremnøder hiin vor Sjels egen Action og Production, saa at vor individuelle
Organismes eller Mikrokosmes Affection for os adskiller Sandsning fra blot
Forestilling.\par}

\begin{center}\large\S~4.\end{center}
\markboth{\S~4}{\S~4}
\begin{center}Philosophie er b) intelligent og rational Erkjenden. --- Erkjendelsens
Grundmomenter i deres Fremtræden som saadanne. --- Erkjendelsen som intelligent.\end{center}

Den individuelle Bevidsthedsleven, der bestemmes ved det Objectives Magt, viser
sig, forsaavidt den træder frem ligeoverfor det Objective, som den, der i sig
selv er fri, medførende en uudtømmelig Muelighed til at virke paa andre Maader,
end den Maade, som den bestemmende Nødvendighed hver Gang medfører.
Nødvendigheden refererer sig altsaa her til en Frihed, og har en saadan til
Gjenstand, idet det netop er det i sig selv Frie (den i sig selv frie Imagination
og Reflexion), som ved Nødvendigheden bindes, og Erkjendelsen faaer ogsaa kun
herved Charakteren af en Erkjenden, idet den kun herved viser sig som en
Gyldighedstillæggelse. [Det Sande træder ogsaa kun i samme Grad frem som det
Gyldige og i sin Gyldighed Erkjendte, som dette Frie, trædende frem som saadant,
nu derpaa viser sig at være Det, der bestemmes og bindes].

Men i den første naturlige, sig af sig selv dannende Erkjenden, ved hvilken vi i
det Hele forholde os som Naturvæsener (cfr.\ min Psych.\ ny Udarb.\ \S~26), træder
Modsætningen imellem det Objective og Subjectiviteten, om den end enkeltviis kan
vise sig, ikke saaledes frem, at den heelt gjør sig gjeldende som saadan.
Erkjendelsen befinder sig da, uagtet alt det i sig selv deri liggende Medierende,
paa et første Umiddelbarhedens Trin, og Erkjendelsens Momenter, om de end vistnok
ere tilstede, ere da saaledes fortabte i den hele Erkjenden, at de ei komme frem
som saadanne. Men da træder ei heller det Sande frem for os i \emph{sin Sandhed},
og det kommer ei til at gjøre sig gjeldende ved sin formelig anerkjendte
Gyldighed.

For at det Sande skal komme til en saadan høiere Hersken, maa først den paa
Subjectivitetens Side liggende individuelle Bevidsthedsleven heelt træde ud i sin
Frihed. Men hertil fører da ogsaa Erkjendelsesbestræbelsen i sin Fremgang, idet
vi --- enten ved anden Interesse, eller ved blot Videdrivt og den sig
herigjennem fremarbeidende høiere Leven --- bevæges til at gaae ud over det
umiddelbart Givne, for at søge Andet og Mere, enten blot a) ved at see os videre
om i den for os aabnede Sphære, hvorved da ei allene Sindshenvendelsen, altsaa
Opmærksomheden, bliver fri, men ogsaa, idet Forventning indtræder, Forestillingen
begynder at finde sig løsbunden; eller ogsaa b) ved at spørge om Andet og Mere,
end det i hiin Sphære Givne, fordi vi ved dette finde os vakte til at søge noget
Saadant, som Det, hvorpaa hiint viser hen, saa at altsaa en Randsagelsesaand gjør
sig gjeldende, og Granskning eller Efterforskning fremkommer, og nu ei allene
Forestillingen, Imaginationen, men ogsaa den deri liggende Forholdsforfølgelse og
Sammenhængsforfølgelse, eller Reflexionen, træder frem som saadan. --- Men det
Nye, som da nu paa denne Maade formeligen søges, søges da tillige som det Gyldige
og til Efterforskningen saaledes Svarende, at hiin Spørgen kan finde sig
tilfredsstillet; hvorved der da tragtes efter et Resultat, ved hvilket den
friblevne, eller i sin Frihed udtraadte, Forestillen og Forfølgen kan
tilstrækkeligen atter vorde bestemmet og bunden, med en Nødvendighed, som da nu
ogsaa søges, og naar den findes, vil komme til at træde frem som saadan.

Men den Nødvendighed, tilligemed den Beherskethed og Indlevethed i det Objectives
Magt, som paa denne Maade indtræder, har da en anden og friere Charakteer, end
hiin første umiddelbare Nødvendighed. Den træder nu frem som Forstands- og
Fornuftnødvendighed, og vi selv vise os nu som Fornuftvæsener, der tillægge
Gyldighed, fordi vi anerkjende ifølge en i Indsigt grundet udtrykkelig
Indrømmelse, Istemmen, Bifaldelse; hvorved da ogsaa Anerkjendelsen eller
Gyldighedstillæggelsen kommer til, trædende ud i sin Frihed, at vise sig som
saadan. Thi om det end maaskee fra først af kun ere nye Gjenstande, af samme Art,
som de paa umiddelbar Maade givne, hvilke saaledes søges, under en blot
eftersporende Iagttagen, saa maae de dog nu ei allene vise sig for os som de, vi
søgte, men vi maae da nu ogsaa forvisses om, at de ere de rette, og at vi nu
objectivere os dem retteligen, saa at intet Andet stiller sig frem for os og af
os statueres, end det i sig selv Gyldige, eller Det, som \emph{bør} statueres.
--- Men Dette faaer da naturligen en saadan Tilbagevirkning paa det tidligere
umiddelbart Antagne, at ogsaa dette kan drages ind under samme Overveielse, saa
at \emph{al} Erkjenden vil kunne blive til Gjenstand for et Spørgsmaal om dens
Gyldighed.

Erkjendelsen vil da paa denne Maade blive en saadan, hvilken vi selv bekræfte hos
os, idet vi bekræfte os i den, saa at nu ogsaa Forvisningen kommer til at vise
sig som saadan, og tilmed som en potentseret Forvisning, eller som en Forvisning,
om hvis Gyldighed vi ere forvissede. --- Men vistnok vil da ogsaa herunder det i
Spørgningen liggende Uvisse kunne komme saaledes frem, at noget Tvivlende,
Dubitativt, kan indtræde, hvilket imidlertid overvindes, naar et Resultat
opnaaes, som bliver tilstrækkelig tilfredsstillende.

% Rettelse (errata leaf, p.14 l.28): print reads ``Tilfredsstillelelse'';
% erratum ``l. Tilfredsstillelse''. Applied silently in ``Tænknings-Tilfredsstillelse''.
Her fremkommer da en egentlig medieret, nemlig ved Tænkning medieret, og paa
Tænkning og Tænknings-Tilfredsstillelse beroende, eller intelligent Erkjenden,
under hvis Dannelse de oprindeligen i den til Erkjendelse henstræbende
Sammenhængs- og Forholdsforfølgelse, d.\ e.\ i Tænkningen selv (i Intelligentsen
i Subjectet), liggende Grundfordringer ere saaledes i Bevægelse, at det er ifølge
deres Tilfredsstillelse, at Indsigten, Bifaldelsen og Anerkjendelsen indtræder.

\begin{center}\large\S~5.\end{center}
\markboth{\S~5}{\S~5}
\begin{center}Fortsættelse. Hvorpaa den intelligente Erkjenden gaaer ud.\end{center}

En saadan intelligent Erkjenden danner sig under al til sit Maal kommende
Forskning idetmindste forsaavidt, at der indsees, a) \emph{at} noget Vist er at
antage og bør at statueres. Men spørge vi, om nu ogsaa Det, b) \emph{hvad} vi
erkjende, eller Erkjendelsens Indhold, constituerer sig for os ved Tænkning,
altsaa ved Forstands- og Fornuftnødvendighed, saa bemærke vi, at ogsaa Dette er
Noget, hvortil Tænkning fører. Vi komme da til at skjelne imellem 1) den sig
selv, ved sig selv --- dog gjennem Tænkningens eller Intelligentsens Medium ---
bekræftende Affectionsnødvendighed med den paa samme hvilende Erkjenden, og 2)
den ved Forstands- og Fornuftbetragtningen i \emph{Henseende til selve det
Antagne} constituerede Nødvendighed med den herpaa hvilende Erkjenden. I første
Tilfælde er Nødvendigheden af hiin første umiddelbare Art, der gik forud for
Randsagelsen, og den gjør sig gjeldende paa samme Maade, kun at vi nu gjennem
Reflexion have funden os berettigede til at give den Rum og Magt. Erkjendelsens
Gyldighed fremtræder da her for os paa Grund af en Forholds- og
Sammenhængsforfølgelse, og Anerkjendelsen eller Gyldighedstillæggelsen har sit
Grundlag i en saadan; men den gjør sig endnu ikke gjeldende i Erkjendelsens
Constitution i Henseende til selve det Erkjendte. [Som naar man ved
Prøvelsesmidler forvisser sig om, at man har observeret et Phænomen rigtig.] I
det andet Tilfælde derimod fastsættes og bestemmes Det, hvad der skal erkjendes,
ved Intelligentsen, saa at Erkjendelsen endog faaer Udseende af at være en
saadan, vi ved egen Tænkning og Bestemmen selv constituere os, uagtet vi dog
ogsaa her ere bestemmede eller determinerede, idet dog ogsaa her
% FLAG: print reads ``Diet'' (capital D, dotted i, e, t) --- semantically odd in
% Danish, not in the errata leaf; transcribed as printed. Verify against copy.
Diet for Det, som maa lede og bestemme denne vor Constitueren, bliver
Grundlaget. Kun den sidste Art af Erkjenden kunne vi egentligen kalde den
intelligente Erkjenden, og den er da en saadan, som deels a) skeer i Kraft af
Sammenhængs og Forholds Forfølgelse, deels b) bestaaer i en Erkjendelse af selve
Sammenhængene og Forholdene samt af det i dem eller for dem til Grund Liggende,
der i dem udgjør det indre Constitutive.

{\small\emph{Anm.~1.} Det vil let sees, at, naar vi under det førstnævnte
Tilfælde henregne den Reflexion, som bestemmer os til at tillægge en fornummen
Affectionsnødvendighed Gyldighed, saa vil derimod Reflexionen over Grundetheden
af denne Gyldighedstillæggelse, naar denne nu betragtes som vor Spørgens
Gjenstand, føre os over i det andet af de to her adskilte Tilfælde.\par}

{\small\emph{Anm.~2.} Fremdeles bemærke vi, at her kunne gives Tilfælde, da en
Erkjendelse af et „hvad" kan, tilligemed sit „at", constituere sig for os paa
intelligent Maade, uden at vi ere trængte ind i Gjenstandens Hvorledes. Saaledes
vil Gyldigheden af vor Tro paa en Styrelse og Gyldigheden af vor Leven i denne
Tro kunne paa intelligent Maade godtgjøre sig for os, og det saa, at selve
Erkjendelsen af en saadan Styrelses Existents constituerer sig for os gjennem
Tænkning, uden at derfor en Indsigt i selve denne Styrelse danner sig.\par}

{\small\emph{Anm.~3.} Angaaende Udtrykket Forstands- og Fornuftnødvendighed,
hvormed her den Nødvendighed, vi og kunne kalde den intellectuelle, er betegnet,
bemærkes, at det dobbelte Udtryk er indtraadt, fordi her nærmest er at tænke paa
Forstanden (\textit{intellectus}), som den, der gaaer ind i den Sammenhængs- og
Forholdsforfølgelse, den baade objective og subjective Tankeforfølgelse, som
herved finder Sted, og som fører til at indsee, og altsaa og at forstaae
(\textit{intelligere}); men at nu dernæst dog, naar Anerkjendelsen finder Sted
som saadan, saa at der statueres, hvad vi see at vi maae, og tilmed \emph{bør}
statuere, Sandsen for det Algyldige i dets Algyldighed er i Bevægelse paa
Besindelsens Maade, altsaa Fornuften.\par}

Men nu kan det Objective dog ogsaa vise sig for os ved --- og tiltale os gjennem
--- det paa velgjørende Maade Fyldestgjørende, som det kan have for os, og kan nu
og i sin Gyldighed gjøre sig gjeldende hos os ad denne Vei. Da komme vi nærmest
til at bevæges saaledes deraf, at vi knytte os til det, hengive os deri og elske
det, hvorved da, naar nu ogsaa en Anerkjendelse heraf udvikler sig, en ganske
anden Art af Erkjendelse indtræder, end Erkjendelsen i Kraft af Indsigt og
anerkjendt Grundethed, nemlig den, jeg andensteds har kaldet den sympathetiske,
og har fremstillet som tredie Hovedart af Erkjenden, i Tredelingen a) den
umiddelbart skuende, b) den intelligente, c) den sympathetiske Erkjenden, til
hvilken sidste deels den æsthetiske, deels Troen og hvad der monne være at stille
sammen med denne, henhører. See min Psych.\ ny Udarb.\ \S~33, P.\ 130;
cfr.\ Maanedskr.\ f.\ Liter.\ 1838, H.\ 6.

{\small\emph{Anm.~5.} Men denne Tredeling maae vi endnu i en anden Henseende her
tage i Betragtning. Vi ere komne til at see, hvorledes Erkjendelsen fra sit
første Hovedtrin, Umiddelbarhedens, eller den første umiddelbare Erkjendens
Trin, gaaer over til at antage en ny Charakteer, ved at træde frem som den ved
Randsagelsesaanden fremkaldte middelbare, men tillige intelligente og derved
høiere Erkjenden. Men nu maa da det Spørgsmaal opstaae, hvorledes denne forholder
sig til hiin sympathetiske Erkjenden, der ogsaa viser sig som en saaledes
medieret, at ogsaa ved samme Subjectiviteten træder frem som saadan ligeoverfor
det Objective. Vi bemærke da, at denne her bliver at henregne med til hiin første
umiddelbare Erkjenden, da den --- om den end i Sammenligning med den umiddelbart
skuende Erkjenden viser sig at være en mere middelbar, der forsaavidt bliver at
sammenstille med den intelligente --- dog, ligesaavel som den umiddelbart
skuende, danner sig forud for al Randsagelse og Prøvelse, og afgiver Noget som
ved den intelligente Erkjenden først skal gjøres til Gjenstand for Spørgsmaal om
dets Gyldighed og dets Indhold. --- Ligeoverfor den intelligente Erkjenden, der
ogsaa, som vi længere hen nærmere ville see, kan kaldes den rationale, finde vi
en anden, til hvilken vi ikke komme gjennem en rational Stræben, men som danner
sig forud for denne, og har sit eget, af det Rationale, hvorved vi mene det i sin
Rationalitet fremtrædende Rationale, uafhængige Grundlag. Denne anden Erkjenden
kunne vi da kalde den extrarationale eller præterrationale (fordi vi ikke have
den ifølge Randsagelse og Regnskabsaflæggelse eller ifølge Grunde, der have
fremstillet sig for os som saadanne, men faae den paa anden Maade). Men til
saadan extrarational Erkjenden hører nu baade den, der allene skylder den
udvortes Sandsning sin Tilvær, samt hvad der bliver at stille sammen hermed
(hvorhen Rums- og Tidsanskuelsen henhører), saa og den hele sympathetiske
Erkjenden, der ogsaa danner og haandhæver sig udenfor hiin Rationalitet. Al
saadan extrarational Erkjenden er ogsaa, i Henseende til sin Opstaaen, ifølge sin
Natur mere beroende paa umiddelbar Paavirkning, Affection, end den intelligente.
Iøvrigt finde vi ogsaa i dennes Sphære Meget, der danner sig paa Umiddelbarhedens
Maade, og maa det overhovedet bemærkes, at Menneskeheden, især i de dannede
Nationer, nu paa Umiddelbarhedens Maade, eller af sig selv, d.\ e.\ uden
udtrykkelig Overveielse og Betænkning, tilegner sig Meget, som man fra først af
kun langsomt har arbeidet sig hen til. Dette er, hvad især de Gamles Musik
angaaer, bleven oplyst af Hr.\ Pastor Bindesbøl i hans Anmældelse af \emph{Dr.}
Bojesens Doctordissertation \textit{de harmonica scientia Græcorum} (1833) i
Maanedskr.\ f.\ Liter.\ 1837, 8de Hefte; cfr.\ Indl.\ til \emph{Dr.} Bojesens
Skoleprogram 1843 \textit{de tonis s.\ harmoniis Græcorum}, Pag.\ 6.\par}

{\small\emph{Anm.~6.} Det her, naar vi see hen til den menneskelige Stræben, som
Stræben, og dennes Gang og Skjæbne i Menneskeheden, naturligen opstaaende
Spørgsmaal: hvorfor dog ikke al Erkjenden, ogsaa den dybere og høiere, giver sig
os paa den simplere, ligefremmere og umiddelbarere Maade, saaledes som dog
Genialitetens Umiddelbarhed viser det enkeltviis at skee i de høiere Sphærer, er
af mig bleven omhandlet i mit Skrivt „om Erkj.\ og Gr." i Anm.\ til \S~4,
Pag.\ 32--36; hvormed det i min Pathologie \S~23, Pag.\ 55--56 Fremsatte kan
sammenholdes. Men jeg vil her vende tilbage til dette Spørgsmaal, for at give en
mere omfattende Besvarelse deraf, end hist er bleven givet.\par}

Tre Ting er der herved at see hen til: Vi bemærke, at vi ikke allene maae a)
langsomt gjennem megen Undersøgelse og Anstrengelse arbeide os hen til den høiere
Erkjenden, men at Veien endog b) gaaer igjennem megen Tvivl og Uvished, ja c)
aabner en saadan Muelighed af Vildfarelse, at denne paa mangen Maade er bleven
til Virkelighed; og vi spørge nu, hvorfor dog det hele Forhold er et saadant.
Hertil er da at svare: a) at det saaledes hører til den individuelle aandelige
Udviklingsgang paa den ene og til det Objectives fuldelige Objectivering paa den
anden Side, saa at baade Personlighedens fuldelige Fremgang paa sin Side og det
Objectives Riges Fremgang paa sin Side maa føre det saaledes med sig; og at, hvad
b) Tvivlen især angaaer, den fuldelige Selvbekræftelse og Sagbekræftelse kun
herigjennem Punct for Punct kan finde Sted; samt at c) Mueligheden til
Vildfarelse ei kunde være Andet, end en ligefrem Følge af det hele Forhold, idet
dette medførte, at Reflexion og Imagination maatte vorde frie og gjøre sig
gjeldende i deres Frihed, for at et Liv i Intelligents og Frihed kunde udvikle og
danne sig, men at et saadant atter var det nødvendige Grundlag for Uddannelsen af
et Rige i Aand og Kjærlighed.

Men tilsidst bliver endnu at bemærke, at, naar der sees hen til Tilstanden i
Menneskeheden, som den i Gjerningen viser sig at være, maa det vistnok tages i
Betragtning, at den almindelige menneskelige Syndighed og den hermed forbundne
Abnormitet eller Udskeielse og Forstyrring i den hele menneskelige Leven ei har
kunnet Andet end gribe hæmmende, forstyrrende, forvanskende ind i den
menneskelige Erkjendelsesstræben og tilmed i Erkjendelsens hele Gang og Art, og
at der hertil maa tages Hensyn, naar vi stille os Virkeligheden for Øie, og finde
den fuld af Uvished, Tvivl og Vildfarelse.

\begin{center}\large\S~6.\end{center}
\markboth{\S~6}{\S~6}
\begin{center}Erkjendelsen som rational, henarbeidende til et sidste Grundlag.\end{center}

Forsaavidt den indtrædende Spørgen og Randsagelse fører til den Prøvelse,
hvorunder vi gjøre os Regnskab for Gyldigheden af den nye Bestemtvorden, som skal
indtræde, viser sig det Bestemmende som Grund, eller som Det, hvorved vi, ifølge
en Regnskabsaflæggelse, finde os bestemte til at antage hvad vi antage
(\textit{ratio}), og forsaavidt et saadant Bestemmende finder Sted Punct for
Punct i den hele nu sig dannende nye Erkjenden, indtræde her paa en fortløbende
Maade bestandigen Grunde (\textit{rationes}), og Erkjendelsen viser sig da som en
Erkjendelse i Kraft af Grunde, og tilmed som den, der gjør sig gjeldende gjennem
den nu som saadan fremtrædende Bevidsthed om dens Grundethed. Den intelligente
Erkjenden viser sig da med Hensyn til baade den Regnskabsaflæggelse og den
Grundenes Overveielse, som herved finder Sted, som en rational. Det i sig selv
Grundede og Gyldige antages da kun, forsaavidt det har viist sig for os i sin
Grundethed og derved i sin Sandhed.

Men naar Grunde søges og forfølges og en fuldelig Forvisning skal søges i
Grundetheden, saa maa der gaaes tilbage til sidste Grunde og et sidste Grundlag,
hvilket saa derpaa bliver at betragte som det første Grundlag, og Grundene som de
første Grunde for den hele Erkjenden, der først herved bliver grundig rational.

{\small\emph{Anm.} Det i sig selv gyldige Sande kommer herved til, idet det nu
træder frem i sin anerkjendte Gyldighed og Sandhed, at herske med en høiere Magt,
idet det kommer til at herske, som Det, vi selv yde vor Hylding. --- Men ogsaa
det modstaaende Subjective bliver herved opløftet til et høiere Trin; thi da
Erkjendelsen her gaaer frem af Randsagelse og Prøvelse og deraf følgende
Samstemming og Bifald, kommer Subjectiviteten til at røre sig med Frihed og med
Harmonie i hiin høiere Nødvendigheds og Behersketheds Sphære. --- Behersket er nu
besjælet, beaandet.\par}

\begin{center}\large\S~7.\end{center}
\markboth{\S~7}{\S~7}
\begin{center}Philosophie er c) videnskabelig Erkjendelse. Om Videnskab og
Videnskabelighed i Almindelighed.\end{center}

Overveie vi nærmere, hvad Grunde ville sige, saa see vi, at, da de bestaae i Det,
som gjør, at, naar noget Vist er givet, noget vist Andet tillige maa være, eller
statueres at være, saa vise de tilbage til en vis Sammenhæng, og denne atter til
en vis Grundsammenhæng, ifølge hvilken det er, at med eller af det Ene det Andet
følger, saa at altsaa Grundes Erkjendelse maa blive Sammenhængs Erkjendelse, og
at Erkjendelse i Kraft af Grunde maa blive Erkjenden i Kraft af Sammenhængs
Forfølgelse. Men saaledes maa da den rationale Erkjenden føre til, at det
Indbegreb af Erkjendelser, som haves, bliver et velsammenhængende, hvori hvert
Enkelt, idet der gjøres Rede for det, tillige sees i sin Indoptagethed i en sig
dannende Totalitet af Erkjendelse. Herved bliver Erkjendelsen videnskabelig, og
Indbegrebet af Erkjendelse bliver til Videnskab i egentlig Forstand. --- Herved
ordner sig en Fleerhed af Erkjendelser, der, forsaavidt de træde frem hver for
sig med sin Bestemthed, udgjøre Videnskabens Detail, saaledes sammen, at der
danner sig en velsammenhængende Heelhed, eller et System, i hvilket imidlertid
selve den Erkjenden, selve den Indseen og Rationalitet, ifølge hvilken alt Enkelt
bestemmes, og det hele System ordnes og dannes, saaledes er Hovedsagen, at det
kommer an paa, at en saadan, som vi nu kalde snart Videnskabens Aand, snart dens
Idee, kommer til at røre sig og leve i hiin Totalitet af Erkjendelser, i hvilken
nu imidlertid dog alt Enkelt maa træde ud hvert for sig, saa at Videnskabens
Detail, ved en intelligent og rational Forfølgelse af selve Sagens Detail,
udfolder sig i et Indbegreb af Begreber og Begrebsbestemmelser, og Systemet viser
sig som en til den objectivt givne Heelhed svarende Heelhed af begrebsmæssig
Erkjenden.

{\small\emph{Anm.} Til den her nævnte i Videnskaben, som saadan, liggende Trehed
af a) Aand eller Idee, b) System og c) Detail, vil endnu, naar der sees paa
Videnskabens saakaldte Behandling, saavel under Studiet, som skal føre til den,
som og ved Meddelelsen til Andre, komme Methoden, der da deels kan være at
betragte i subjectivlogisk, deels i didaktisk Henseende. Men herved er der da ei
længere Tale om Sagen selv, men kun om vor Erkjendelses- og
Fremstillingsbestræbelse. Angaaende saavel hiin objective Trehed, som dette
Subjective, henviser jeg iøvrigt til min Bog om Erkj.\ og Granskn., hvor der
ogsaa (i \S~17) er Tale om det indre organiske Vexelvirkningsforhold imellem
Idee, System og Detail. Kun ved dette Vexelvirkningsforhold vil jeg her dvæle
noget, og bemærker herom Følgende:\par}

Vi befinde os midt i en stor Totalitet af Gjenstande og Forhold, og det ei allene
i udvortes Henseende, men ogsaa i Henseende til Intelligentsens og de høiere
Forkyndelsers og Virkningers Rige inden i os. Vort Liv danner sig herunder paa en
sporadisk Maade, ud fra utallige Puncter, der gjøre sig gjeldende, hvert Punct
paa sin Maade, og hertil svarer da et heelt Indbegreb af Iagttagelser eller
Bemærkninger, hvilke da hver træde frem med sin, virkelige eller dog
formeentlige, Gyldighed. Saadanne kunne ogsaa fremdeles med Flid anstilles paa
utallige Puncter, idet da hver Gang Det iagttages og opfattes, hvad der hver Gang
træder frem. Først naar nu et Indbegreb af saadanne er samlet, kan den
Forfølgelsesaand begynde at virke, som skal føre til den velsammenhængende
Totalerkjendelse, til System og systematisk Indsigt. Her vil da den isolerede
Iagttagen af alt det Enkelte i dets Enkelthed bestandig gjøre sin Ret gjeldende,
saa at den systematiske Erkjenden maa blive, som den detaillerede Bemærken fører
den med sig, men omvendt vil den systematiske Indsigt atter have en bestemmende,
opklarende, rensende Indflydelse paa den detaillerede Bemærken. Heelheden vil
saaledes i og for vor Erkjenden bestemmes ved Detaillet, og Detailerkjendelsen,
samt Maaden at see og stille Detaillet paa, atter bestemmes ved
Heelheds-Erkjendelsen.

% Rettelse (errata leaf, p.24 l.13): print reads ``snaledes''; erratum ``l. saaledes''.
% Applied silently below (``det Enkelte tages saaledes'').
Men i Heelheden, som det Mangfoldiges Sammenfattelse eller Sammentræden til
Eenhed, maa der dog nu fremdeles aabenbare sig Det, i Kraft af hvilket baade det
Enkelte tages saaledes, som det tages, idet det paa en bestemt Maade indoptages i
Heelheden og indordnes i denne, saa og Heelheden organiseres saaledes, som den nu
herunder udføres. Og saaledes komme vi til et dybere til Grund Liggende: den
ligesom allestedsnærværende Grunderkjendelse, som hæver sig frem under den hele
Iagttagen og Organiseren, og nu baade udgjør det Grundbestemmende og
Grundveiledende, og tillige dog selv først uddanner og antager bestemt Skikkelse
under hiin Stræben. Denne kalde vi nu snart Ideen, snart Aanden i Videnskaben.

\begin{center}\rule{0.18\linewidth}{0.4pt}\end{center}

%% ============================================================
%%  PROPÆDEUTISK BEGREB OM PHILOSOPHIE OG SAMMES EXPOSITION.
%%  (§§ 8--14, printed pp. 25--58; PDF 36--69)
%% ============================================================
\grouphead{Propædeutisk Begreb om Philosophie og sammes Exposition.}

\begin{center}\large\S~8.\end{center}
\markboth{\S~8}{\S~8}
\begin{center}Philosophie gaaer d) ud paa en altomfattende Grunderkjendelse.
Subjectiv Hovedopgave.\end{center}

Et Indbegreb af Erkjendelser kan være intelligent og rational, samt tillige saa
aldeles videnskabeligt, at det fører til en egentlig Videnskabs Uddannelse, uden
derfor at være Philosophie, ja endog kun at komme i Berøring med Philosophien,
efter det gjeldende Begreb om samme, hvilket vi her (ifølge det i \S~1 Bemærkede)
have at holde os til. Herpaa finde vi baade i Naturvidenskaberne og i Mathematiken
talende Beviser. Der spørges da, hvad der er det Eiendommelige ved den Videnskab,
som kaldes Philosophie.

Vi bemærke da allerførst, at Philosophien skal føre os tilbage til de sidste
Grunde. Men herved opstaaer nu det Spørgsmaal, hvad der herved menes, og da bliver
Følgende at tage i Betragtning:

Det var Erkjendelsen, der, under den rationale Stræben, nærmest og først viiste
sig som den, der gik over til at vorde en saadan, der kom til at finde Sted i
Kraft af Grunde; det var Erkjendelsen, som skulde vise sig for os i sin Grundethed
og Gyldighed og derved i sin Sandhed. Forsaavidt vare altsaa de Grunde, hvorom
der var Tale Erkjendelsesgrunde (\textit{rationes cognoscendi}). Her spørges da
bestandigen: hvad skulle vi antage og statuere, og med hvad Ret antage vi hver
Gang Det, vi antage og ifølge denne Antagen erkjende; men Bestræbelsen gaaer da
her hen til et sidste Grundlag.

Herved føres vi nu da nærmest og først til at see hen til Erkjendelsernes første
Udspring eller Kilde, hvorved menes Det, hvorved de gives os, eller som udøver
den bestemmende Indflydelse paa os. Men her fremtræder da i forskjellige
Erkjendelsens Sphærer en forskjellig Art af umiddelbart Bestemmethed, hvilken da
i hver Sphære for sig giver noget Første, hvorfra vi, ved at forfølge de deri sig
visende Forhold, ledes hen til Mere, hvilket da atter enten, naar det nu træder
frem for os, a) gjør sig gjeldende ved samme Art af umiddelbart Bestemmende, som
hiint Første, saa at det Medierende kun bliver et Henledende, eller ogsaa b) viser
sig som Det, der paa intelligent Maade antages paa Grund af hiint Første, eller
antages som Noget, der indsees at maatte være en Følge af hiint, uagtet Dette ei
paa samme umiddelbare Maade staaer for os, som hiint. I første Tilfælde kan
Prøvelsen kun gaae ud paa, om vi ogsaa virkelig vorde bestemmede (afficerede)
saaledes, som det fremstiller sig for os (om vi ogsaa virkelig see, høre
o.\,s.\,v.\ Det og Det saa og saaledes); i det andet Tilfælde gaaer Prøvelsen
saaledes ud paa Det, hvad der skal antages, at det skal sees at svare til og
stemme sammen med det dertil henførende Begrundende, saa at her bliver Tale om
intelligent Tilfredsstillelse (f.\ Ex.\ ved at bestemme den Forestilling, man skal
danne sig om Planetsystemet). See ovenfor \S~5.

Men af en ganske anden Natur bliver Forfølgelsen og Spørgsmaalet om vor
Erkjendens Gyldighed, naar det Dubitative, som her kan fremkomme, og det
Problematiske (især under store Collisioner, idet fra een Side Eet, fra en anden
et Andet, synes at ville gjøre sig gjeldende) fører til at spørge om al menneskelig
Erkjendens Gyldighed, tilmed al Forvisnings Grundethed, overhovedet, saa at altsaa
Paalideligheden af alle hine relativt første Kilder og Gyldigheden af deres
Auctoritet, tilmed den objective Sandhed af alt herigjennem Givet, gjøres til
Gjenstand for et Spørgsmaal og en Overveielse, der da bliver saaledes altomfattende,
at et sidste Grundlag for al Erkjenden, som Erkjenden, eller for al Antagen og
Gyldighedstillæggelse overhovedet, og en i denne Henseende altomfattende
Grunderkjendelse fordres. Men man seer let, at der da ei længere er Tale om
Erkjendelsesgrunde, men om al Erkjendelses egen Grundethed og Gyldighed, saa at
der her endog bliver Tale om Andet og Mere, end om Det, vi kunne kalde
\textit{universam cognoscendi rationem}, naar vi nemlig herved mene al Erkjendens
Grundnatur og Grundbeskaffenhed. Herved træde vi da over paa Philosophiens Gebeet,
ja baade Sagens egen Natur saa og et Henblik til Philosophiens hele Historie maa
føre os til at opstille det her opstaaende Spørgsmaal som den ene af Philosophiens
Hovedopgaver, hvilken vi da ville kalde Philosophiens subjective Hovedopgave. Men
herom komme vi til at tale længere hen, hvor det da vil vise sig, at denne Opgave
optages i en større Opgave.

\begin{center}\large\S~9.\end{center}
\markboth{\S~9}{\S~9}
\begin{center}Objectiv Hovedopgave. --- Philosophie gaaer e) ud paa Tilværelsens
sidste Grunde.\end{center}

Den hele Rationalitet i vor Erkjenden gaaer nemlig ud paa Mere, end blot
Erkjendelsernes Grundethed og Gyldighed. Den gaaer ud paa selve Gjenstandenes
Grundforhold, og altsaa tillige ud paa disses egne Grunde, idet vi stræbe efter,
at Det, hvori hvert Punct i selve Tilværelsen er grundet, maa vise sig for os. Thi
vor hele Erkjendelsesbestræbelse gaaer jo ud paa, at vi i og gjennem Erkjendelse
kunne fuldeligen tilegne os og indleve os i Objectiviteten; men da maae vi gaae
ind i Sammenhængen og Forholdene i alt det Givne, for at Alt heri kan vise sig for
os, hvert i sin Betydning i Totaliteten. Men saaledes komme vi da til, ei blot at
kræve os selv til Regnskab med Hensyn til vor Opfattelse af det Objective, men
komme ogsaa til at kræve selve det Objective paa ethvert Punct, eller kræve selve
Tingene til Regnskab, for at de kunne ligesom gjøre os Rede for, hvori de selv ere
grundede, og hvorledes det staaer sig med deres Grundforhold. --- Men idet vi nu
ogsaa herved tragte efter at gaae til sidste Grunde og Grundforhold, komme vi,
idet hele Tilværelsen gjør sig gjeldende for os som den, der danner en eneste
Heelhed, til at tragte efter en hertil svarende altomfattende Grunderkjendelse,
ikke som et Aggregat af flere, men med indre Eenhed og Heelhed i sig.

Og herved træde vi da atter ind paa Philosophiens Gebeet og møde dens anden
Hovedopgave, hvilken vi ville kalde den objective.

\begin{center}\large\S~10.\end{center}
\markboth{\S~10}{\S~10}
\begin{center}Propædeutisk Definition paa Philosophie.\end{center}

Om Philosophien overhovedet kunne vi da sige, at den fremkommer, naar
Bestræbelsen efter a) at gaae til de sidste Grunde, tillige gaaer ud paa b) at
blive altomfattende og det saaledes, at den vil føre Alt tilbage til c) eet
universelt Grundlag, saa at een altomfattende sidste Grunderkjendelse kan danne
sig. Og hermed er da den første propædeutiske Definition paa Philosophie given.

Men naar vi nu tilføie, at den, som vi have seet, og som ogsaa Philosophiens
Historie lærer, fra først af har en dobbelt Hovedopgave, da baade Spørgsmaalet om
al vor Erkjendens sidste Grundlag og Grundethed, Realitet og Objectivitet, Sandhed
og Vished, saa og Spørgsmaalet om selve Tingenes sidste Grunde og Grundforhold,
træder frem som Hoved- og Grundspørgsmaal, komme vi, ved at overveie begge
Hovedopgavers Forhold til hinanden, til at indsee, at den subjective Opgave bliver
at henføre under den objective. Thi al vor Erkjendens sidste Grund og
Grundbeskaffenhed (\textit{universa ratio cognoscendi}) maa dog beroe paa vor hele
Stilling med Hensyn til det Objective, vor hele Indlemmethed i og Grundforhold til
den sig for os fremstillende Totalitet, saa at, ligesom vi dog kun finde os selv
som Lemmer og Dele af en stor Heelhed, saa maae vi og finde vor Bevidsthed og
altsaa ogsaa vor Erkjenden saaledes medindgaaende i denne Heelhed, at vort hele
Bevidsthedsliv med al vor Erkjendelse i dennes Heelhed kun kan blive een med
blandt de Hovedgjenstande, om hvis sidste Grundforhold og Grundvæsen der bliver
at handle.

Saaledes vil da vor Definition nærmere bestemme sig saaledes, at vi betragte
Philosophien overhovedet som den, hvis Opgave det er paa en altomfattende Maade
at trænge hen til hele Tilværelsens eller alle Tings sidste Grunde og Grundforhold,
for i Kraft af den herved vundne Grunderkjendelse at opklare Alt. --- Philosophering
er da den tænkende Betragtning eller Tankeforfølgelse, hvori deels en Stræben efter
at naae hen til en saadan sidste, eller som den i en anden Henseende ogsaa kan
kaldes, første Grunderkjendelse arbeider sig frem, deels en saadan Grunderkjendelse,
forsaavidt den har dannet sig, videre søger at gjøre sig gjeldende.

{\small\emph{Anm.~1.} Naar vi erkjende alle Ting i deres sidste Grunde, maae vi
ogsaa komme til at erkjende selve vor Erkjenden, der hører med til disse alle Ting,
i dens sidste Grundethed. I den totale Tingenes Orden, i hvilken vi røre os, finde
vi blandt Andet erkjendende Væsener; disses Tilvær med den hele deri fremtrædende
Stræben efter Erkjendelse maa da være grundet i denne Tingenes Orden, og maa i
Henseende til sin Natur og Tilværelsesmaade og i Henseende til sin Realitet beroe
herpaa. --- Dette vil ogsaa sees, naar den subjective Hovedopgave forfølges. Selv
den subjectivt-idealistiske Løsning deraf, ifølge hvilken det antages, at den sig
for os fremstillende Tilvær ingen egentlig Objectivitet har, men kun har Realitet,
Gyldighed og Tilvær i og for vor Bevidsthed, eller vor Viden af og med os selv,
dreier sig dog, idet den vil bestemme vor Forestillens Realitet, om det
Spørgsmaal, hvad vi skulle tage den hele os givne Objectivitet for at være; og
naar den lærer os, at vi kun leve i en sig for os fremstillende Verden, om hvis
Realitet, uafhængigt af denne dens Sigfremstilling for os, vi aldrig kunne tale,
da vi aldrig kunne have med Andet at gjøre, end det sig for os Fremstillende,
bestemmer den dog derved baade vor Tilvær og Tilværelsesmaaden af Det, vi kalde
det Objective, eller kalde Tingene ligeoverfor os; og den gjør det da saaledes, at
vi da dog ogsaa finde os ligeoverfor dem, og altsaa maae regne os selv med til
Tingenes Totalitet; thi Ting vil her overhovedet kun sige Noget, hvorom der er
Tale, og som kan være Gjenstand for en Antagen og Dømmen. Men Erkjendelsen finde
vi dog nu atter, som noget vist Bestemt midt inde i hiin Totalitet; og er denne
Totalitet kun et System af nødvendige Forestillinger, saa finder sig vor
Erkjendelsesbestræbelse dog nu indkjæbet med som et Led i dette System af
nødvendige Forestillinger. Men især sees den subjective Hovedopgaves
Indbefattethed under den objective, naar vi komme til den Løsning af samme, som
ovenfor (\S~3, Anm.~3) anmærkningsviis er bleven fremsat, eller naar vi endog blot
gaae den hist fulgte Vei, og denne Hovedopgave har derfor ogsaa, under en anden
Form, i Philosophiens historiske Udviklingsgang kunnet fremkomme og maattet
fremkomme saaledes, at Sagen blev aldeles optaget under de objective Opgavers
Forfølgelse, nemlig som Spørgsmaal om Forhold imellem Sjæl og Legeme og dermed
forbundet Forhold imellem Individualitet og Omverden; hvorved da baade Spørgsmaal
og Løsning har dannet sig saaledes, at Opgaven ogsaa kunde og maatte udstrække sig
til Planterne, ja til Alt, hvad der, som Led i Verden, staaer under en Omverdens
bestemmende Indflydelser, og dog virker udaf sin eiendommelige Natur.\par}

% Rettelse (errata leaf, p.31 l.27, NB): print reads ``paa en altomfattende
% opstille''; the NB erratum inserts ``Maade'' (``altomfattende Maade''). Applied.
% Also p.31 l.26 (non-NB): print reads ``opjective''; read ``objective'' — applied
% silently below (``dens objective Opgave'').
{\small\emph{Anm.~2.} Her opstaaer det Spørgsmaal, hvorvidt en saadan Lære, som
Lockes, der kun gaaer ud paa at oplyse al Erkjendens Grundbeskaffenhed og Udspring
eller Kilde, som den under den nøiagtige psychiske Iagttagelse forefindes at være,
skal kunne, saavel i og for sig, som med Hensyn til sin Virkning, henregnes ---
ikke til explicativ Philosophie allene, hvorunder den ligefrem henhører, --- men
og til egentlig eller speculativ Philosophie, uanseet forresten hvad Værd og
Gyldighed vi ville indrømme den, hvorpaa det her ei kommer an. --- Den gaaer i
Grunden kun ud paa at oplyse hvad jeg har kaldet \textit{universam cognoscendi
rationem}; men den er traadt frem som --- eller dog bleven gjort gjeldende som ---
den, hvori ei allene Philosophiens hele subjective Hovedopgave skulde finde sin
Afgjørelse og Beroligelse, men ogsaa Grundvolden for Spørgsmaalet om alle Tings
sidste Grunde og Grundforhold saaledes skulde vorde oplyst, at der i denne
Henseende ei skulde kunne henvises til Andet end Empirien, saa at altsaa speculativ
Philosophie i objectiv Henseende dermed skulde være afviist, og saaledes altsaa
den objective Hovedopgave, som speculativ Opgave, decideret. Forsaavidt maa da en
saadan Philosophie betragtes som den, der har havt at gjøre med baade Philosophiens
subjective og dens objective Opgave. Den har villet paa en altomfattende Maade
opstille en Grunderkjendelse angaaende al Erkjendens sidste Grunde, og angaaende
Veien til alle Tings sidste Grundforholds Udfindelse og Bestemmelse.\par}

{\small\emph{Anm.~3.} Angaaende en Skepticisme, som Humes, er noget Lignende at
bemærke; men den græske Skepticismus maa nu især, som den, der er udsprungen af de
grundphilosophiske Bevægelsers hele Gang og bærer Præget heraf, erkjendes for
Philosophie i egentlig Forstand, forsaavidt her nemlig, som hist, kun er Tale om
det propædeutiske for den speculative Philosophies Historie gyldige Begreb om
Philosophie, og forsaavidt tillige en Totalitet af Philosophering, uden at føre
eller skulle føre til en philosophisk Videnskab, men kun til Uddannelsen af en
philosophisk Anskuelsesviis og Maade at tage Livet paa, kan kaldes Philosophie.
--- Men det sees da nu og, at det her givne Begreb om Philosophie passer sig
aldeles paa hiin Skepticismus; thi ikke, at der paa det empiriske Standpunct Intet
gives, som med Tilforladelighed kan antages, men at, naar der gaaes til de sidste
Grunde, Intet kan opstilles som uanfægteligt vist, er den skeptiske Lære, kun at
den selv naturligviis ei kunde træde op som en Lære, da den naturligviis maatte
indrømme, at ogsaa den maatte blive Gjenstand for sin Skepsis og altsaa i sidste
Instants ei turde træde op som en egentlig Grunderkjendelse, hvorfor den ogsaa kun
skulde gjelde for en sidste Grundberoligelse, der skulde gaae ud paa at bevirke
Sindets ἀταραξία.\par}

{\small\emph{Anm.~4.} Naar Kant, for at virke til en lignende Beroligelse, har
søgt at godtgjøre, at der kun kan være Tale om en blot subjectiv Realitet og
Gyldighed af al vor Erkjenden, men at en saadan nu ogsaa staaer saa fast, som det
for os kan være fornødent, hvorimod al egentlig speculativ Stræben vilde være et
Forsøg paa en Transscendents ud over al vor Fornufts Formaaen, saa har han
dog herved ei allene behandlet den subjective Hovedopgave, men har ogsaa søgt at
afgjøre den objective, nemlig ved at afskjære eller afvise den. --- Fichte derimod,
i den Periode, da han erklærede sin Videnskabslære for kun at være gjennemført
Kantianisme, og kunde erklære den herfor, gik dog netop ved denne Gjennemførelse
forsaavidt ud over Kant, som han i Grunden tillod sig hiin Transscendents, idet
han tog Sagen speculativt; thi han søgte speculativt at opklare og retfærdiggjøre
den Kantiske Lære ved at føre den tilbage til et dybere Grundlag.\par}

% Rettelse (errata leaf, p.33 l.4, NB): print reads ``Begrebets Philosophie'';
% the NB erratum reads ``Begrebet''. Applied below (``af Begrebet Philosophie'').
{\small\emph{Anm.~5.} Det maa her overhovedet bemærkes, at vor Bestemmelse af
Begrebet Philosophie lader det uafgjort, om der ere saadanne Tingenes sidste
Grunde, som af os kunne paa rational Maade erkjendes. Vor Begrebsbestemmelse
bliver derfor ogsaa gyldig for Dem, som af religiøse Grunde og med Hensyn til den
positive Aabenbaring forkaste al den Philosophie, der vil vove sig ind i det
Guddommeliges Region; thi de gjøre det, fordi de nægte, at den menneskelige
Fornuft, idetmindste i Menneskets faldne Tilstand, er istand til at kunne trænge
hen til de høieste Sandheder, hvilket da er det Samme, som at trænge hen til
Tilværelsens sidste Grundforhold. --- Det maa iøvrigt overhovedet bemærkes, at vi
ikke have opstillet Philosophien som en Lære om Tingenes sidste Grunde, men kun
som en Grunderkjendelse, der danner sig under en Stræben efter at komme hen til de
sidste Grunde, altsaa en Lære med Hensyn til dem, eller angaaende dem.\par}

\begin{center}\large\S~11.\end{center}
\markboth{\S~11}{\S~11}
\begin{center}Philosophien træder det ellers Givne imøde. --- Overblik over
Indbegrebet af alt dette.\end{center}

Som den, der skal gaae ud paa at trænge hen til Tilværelsens eller alle Tings
sidste Grunde og Grundforhold, forudsætter Philosophien, at Tilværelsen med Alt,
hvad den indeholder, har gjort sig gjeldende hos os, og er bleven opfattet og
kjendes, og nu fremstiller sig som Det, Philosophien træder hen til. Denne danner
sig da saaledes ligeoverfor og i en vis Modsætning baade til den hele givne
Tilværelse, med hvad denne indeholder, saa og til det hele Indbegreb af al den
Erkjenden, hvorigjennem alt Det, som nu i Philosophien skal finde sin sidste
Opklaring, forud gjør sig gjeldende hos os. Alt, hvad nu Philosophien saaledes
tænkes at træde hen til, kalde vi --- uden dog at ville nægte, at ogsaa
Philosophien kun kan danne sig ved at gives os --- Indbegrebet af det Givne,
hvilket forresten deels gives os paa empirisk, deels paa apriorisk Maade, og i
sidste Tilfælde atter deels paa sympathetisk, deels paa intelligent Maade.

Et Overblik over alt dette Givne vil lære os, at vi deri have at skjelne imellem
følgende Sphærer:

1) Den udvortes Erfarings eller Empiries (Naturforskningens og Menneskehistoriens)
Gebeet, i hvilket Erkjendelserne danne sig ad den explicative Forfølgelses Vei, og
Erkjendelsens første Grundlag bliver Iagttagelse ved Hjelp af de udvortes Sandser,
enten under en umiddelbar Vaerbliven eller Opleven, eller under Opfattelsen af Det,
hvad Andre saaledes have iagttaget og ladet komme til os gjennem saakaldet
Overleverelse (Tradition); men Hovedsagen bliver Forstandens Indtrængen i det
saaledes Givne; som da egentlig Erfaringsvidenskab kun kan bestaae i Det, hvad der
af det umiddelbar Sandsede eller Oplevede samt af det Overleverede uddrages, eller
fremhæver sig, med blivende Gyldighed og til blivende Benyttelse.

2) Den selvstændige aprioriske Constructions eller Operations Gebeet, hvori vi selv
construere, constituere, frembringe Gjenstandene ved en Virksomhed, der ei er
eftergjørende, men gaaer ud paa at realisere forud fattede Tanker, hvilke da danne
sig \textit{a priori}. Det bliver da her dette Aprioriske (Forstandens egne indre
Tilsagn), som bliver Studiets Gjenstand. Men det Aprioriske, vi her have med at
gjøre, gaaer endnu kun ud paa at constituere endelige Gjenstande og at reflectere
over de herved fremtrædende Forhold. --- I de mathematiske Videnskaber, baade den
rene Mathematik og de mathematiskt-mechaniske Videnskaber, finde vi det største og
meest talende Exempel paa et saadant Gebeet; men Exempel kan og findes andensteds,
f.\ Ex.\ i Versificationen og i flere af den poetisk-æsthetiske Konsts Sphærer,
ligeledes i visse Partier af saadanne Videnskaber, som ellers ikke høre herhen,
hvorpaa Syllogistiken, den grammatikalske Syntaktik i visse Henseender, m.\ fl.,
kunne afgive Exempler.

3) Den indvortes sjælelige eller psychiske Iagttagelses Gebeet, hvilket vi komme
til at kalde den explicative Philosophies Gebeet.

Philosophiens Historie maa nemlig bringe os til at adskille to Slags Philosophie,
idet vi fra a) den forhen omhandlede egentlige Philosophie, hvilken vi kalde den
speculative, men ogsaa kunne kalde den constitutive, adskille b) et stort Indbegreb
af Videnskaber, som ogsaa ere blevne kaldede Philosophie, og gaae under dette Navn,
uagtet de kun ere explicative, eller forfølge hver især sit visse bestemte Givne i
dets Givethed, uden at gaae tilbage til alle Tings sidste Grunde, og saaledes
henhøre under de explicative Videnskabers Gebeet, ligesom de forhen under Nr.\ 1
nævnede.

Til dette den indvortes sjælelige Iagttagelses eller den explicative Philosophies
Gebeet hører:

% Rettelse (errata leaf, p.35 l.20): print reads ``for vor Iagttagen og Forfølgen,
% ere''; the erratum ``udslettes Kommaet'' deletes the comma after ``Forfølgen''.
% Applied below.
a) Den indvortes Erfarings Gebeet, Psychologien, hvori Gjenstandene for vor
Iagttagen og Forfølgen ere vore indvortes Sjæletilstande med Alt, hvad der rører
sig i dem, forsaavidt det nemlig \textit{de facto} gjør sig gjeldende; (saa at der
ei er Tale om noget Bør, Noget, som skal paaagtes som det Algyldige, om end derimod
vistnok Sandsen for det Algyldige og Menneskets Forhold til samme, saaledes som det
\textit{de facto} er givet, her maa være at tage i Betragtning). Den Deel af de
herhen hørende Betragtninger, som angaae al Erkjendelses Maade at danne sig paa, er
især bleven kaldet Ideologie.

b) Den blotte Reflexionsphilosophies Gebeet, hvori det Aprioriske, som ligger til
Grund for vor hele Tænkning og for den herved fremkaldte Forfølgelse af Forholdene
i Verden, --- og saaledes altsaa ligger til Grund for Empirien, som styrende og
ledende den hele empiriske Opfattelse og Forfølgelse, --- iagttages og expliceres i
sin Almeengyldighed eller Universalitet (altsaa tillige som det for os Gyldige). Her
bliver da i det Factiske studeret det deri fremtrædende Almeengyldige. Herhen hører
den explicative Logik og Ontologie.

c) Den explicative Ideephilosophies Gebeet, hvori det Uendeliges og Eviges indre
Forkyndelser, forsaavidt de træde frem som saadanne, og tillige vise sig, som de
høiere, paa hvilke det deels beroer, hvad vi selv skulle besjæles og beherskes af,
deels, hvorefter Livets Goder retteligen og grundigen ere at vurdere, blive
Gjenstande for Explication og explicativ Opklaring. Herhen høre den explicative
Religionsphilosophie, Moralphilosophie og Æsthetik.

Men see vi nu tilbage til den speculative Philosophie, der nu skal træde hen til
alle disse Gebeter, saa bemærke vi, at, da den dog ogsaa har gjort sig gjeldende i
Verden, og saaledes forefindes i Verden, maa vor speculative Philosophering dog
ogsaa træde i Forhold til det herved Givne, og saaledes komme vi da til, endnu at
tilføie:

4) selve den speculative Philosopherings Gebeet, forsaavidt det er os givet paa
historisk Maade, eller Philosophiens Historie.

{\small\emph{Anm.~1.} Det har sin Vanskelighed at adskille tilstrækkeligen
Mathematikens hele Gebeet fra det øvrige Aprioriske, saa at en væsentlig Forskjel
kan komme tilsyne imellem Mathematik og explicativ Philosophie. Hvad nemlig
Mathematikens Forskjellighed fra den speculative Philosophie angaaer, da sees den
let. Mathematiken gaaer kun tilbage til sine egne sidste Grunde, sit eget sidste
Grundlag, og er paa ethvert Punct i Stand til at føre Alt tilbage dertil. Men den
gaaer ikke tilbage til, eller søger at gaae tilbage til, alle Tings sidste Grunde,
ei heller til Spørgsmaalet om al Erkjendens sidste Grund og Grundethed. Den har
overhovedet aldeles Intet med selve Tilværelsens Opgaver at gjøre; den bevæger sig
kun i et Muelighedernes Rige (Noget, den iøvrigt har tilfælleds med en blot
explicativ formel Logik), og er den af alle Videnskaber, som mindst af alle kommer
i Berøring med det Speculative. --- Men hvi henregnes den ikke under den explicative
Philosophie? Betragte vi især --- ikke Geometrien, hvilken Kant, som Schmidten
bemærker i sit Skrivt „om Mathematikens Væsen" Pag.\ 30, allene synes at have havt
for Øie, hvor han vil bestemme Forskjellen imellem Philosophien og Mathematiken ---
men Arithmetik og Algebra, som Functionslære, saa maae vi erkjende, at der kan
spørges, hvorfor Mathematiken, forsaavidt den oplyser Calculeringens Natur og
Væsen, ikke ligesaavel er at kalde philosophisk, som den explicative Logik, der
oplyser Tænkningens Natur og Væsen; især da Calcul baade subjectivt deels ligger
til Grund, deels kan komme til at ligge til Grund, overalt i enhver nøiagtig
Opfattelse af det Givne og indtrængende Erkjendelse deri; saa og objectivt er at
betragte som til Grund liggende for selve Tilværelsens Hele, ligesom det Logiske har
en saadan objectiv Realitet og Tilværelse som Grundlag for Tilværelsen; saa at, hvad
enten explicativ Logik regnes til Philosophie, fordi Det, den sætter ud fra
hinanden, hører til det almene Fundamentelle, eller fordi den explicerer det indre
Aprioriske i os, saa maatte der menes at være Grund til ogsaa at henregne
Mathematiken dertil.\par}

{\small Herved er da at bemærke. Det er os givet, at Mathematik ikke henregnes til
Philosophie, den simple explicative Logik og Ontologie derimod, forudsat at den
tages paa den rette grundige Maade, henregnes dertil. Denne Talebrug, samt den
derigjennem udtalte Anskuelse, er det, vi skulle oplyse ved at udfinde, hvad det er,
der gjør, at den kan have dannet sig. Altsaa: vi skulle opsøge den Forskjel, der her
maatte vise sig. Og saa kunde vi da først bemærke, at Logik og Ontologie dog altid
gaae ud paa det mere Fundamentelle, det sidste Fundamentelle; thi, hvad de
explicere, ligger til Grund for Mathematiken, lige saa vel som for al anden
Videnskab. Det er et sidste til Grund Liggendes ligefremme Fremtræden, som her
expliceres. Imidlertid vilde dog denne Forskjel ei kunne være nok, da dog
Psychologien heelt regnes til den explicative Philosophie, uagtet den ei heelt gaaer
ud paa noget almeent Fundamentelt. Det, der maa blive Hovedsagen, er, at den
Opfattelse og Iagttagelse og den Forfølgelse af det Iagttagne, som i de anførte den
explicative Philosophies Gebeter finder Sted, er at kalde en psychisk, i en
Forstand, hvori den i Mathematiken indtrædende ei kan kaldes saa. Thi vel iagttages
ogsaa i denne, hvad det Aprioriske i os har at sige os med Hensyn til de
forekommende Forhold, som skulle forfølges, og de Opgaver, som skulle løses; men
selve disse Forhold, selve Det overhovedet, som iagttages, altsaa Gjenstandene for
Betragtningen, kunne ei siges at være af psychisk Natur (thi om Mathematikens
muelige Anvendelse i Psychologien er her ikke Tale). Derimod iagttages i Logik og
Ontologie deels \emph{Bevidsthedens egne Kjendsgjerninger}, deels det
\emph{Bevidsthedslivet selv constituerende Aprioriske}; hvorimod, naar i
% FLAG: print reads ``dettte'' (triple t) at ``hvori dettte Aprioriske'' below ---
% apparent printer's typo; NOT in the errata leaf; kept as printed. Verify vs copy.
Mathematiken Parablens Formel udfindes, eller nogensomhelst Function bestemmes og
forfølges, Dette vel kan siges at være Noget, som vi lære allene af det Aprioriske i
os, men ikke Noget, hvori dettte Aprioriske selv bliver iagttaget, ikke heller
Noget, som constituerer vor egen Sjæletilvær. Vilde Nogen derimod eftervise det
Mathematiske som Grundlag i vor Reflexion over mange Forhold i Verden, og tilmed i
vor sjælelige Leven, da blev Dette vistnok explicativ Philosophie, men hørte da ei
til Mathematiken, men til Psychologien.\par}

{\small\emph{Anm.~2.} Naar i og gjennem det Factiske det Algyldige skal studeres,
saa bliver Opfattelsen af det Factiske af væsentlig Betydning i og for Læren om det
Algyldige. Dette gjelder især med Hensyn til den explicative Ideephilosophie, hvor
derfor --- uagtet den heelt igjennem maa komme i Stand ved egen Iagttagelse, saa at
Andres Iagttagelser kun ere at tage som de, der skulle tjene hertil, ---
Videnskabernes Historie bliver af langt anden og større Vigtighed for selve
Videnskaberne, end den, som f.\ Ex.\ Mathematikens Historie kan have for
Mathematikerne. Men at gaae saavidt, at man lader den hele Religionsphilosophie
bestaae i Læren om Religionens historiske Udviklingsgang, saaledes som det i vor Tid
er skeet, er ei at tage Sagen retteligen.\par}

% Rettelse (errata leaf, p.39 l.17--18, NB): print reads ``Theologie''; the NB
% erratum reads ``Teleologie'' (which fits ``det Nyttige'' in the parallel list; a
% reader's pencil correction sits beside the word in the scan margin). Applied below.
{\small\emph{Anm.~3.} Det har fremdeles ogsaa sin Vanskelighed at adskille
tilstrækkeligen imellem Reflexionsphilosophie og explicativ Ideephilosophie, da i
begge det Algyldige træder frem, som saadant, saa at et Bør ligesaavel gjør sig
gjeldende i Logiken, som i Moralen, hvorfor ogsaa Treschow henregnede Logiken til
den praktiske Philosophie (hvilken han i sit Skrivt „om Philos.\ Natur og Dele",
inddeelte i Logik, Moralphilosophie, Æsthetik, og Teleologie, efter Formaalene: det
Sande, Gode, Skjønne, og det Nyttige). Men naar man bemærker, at den Paamindelse at
respectere de logiske Fordringer for Sandhedens Skyld, altsaa af Sandhedskjærlighed,
er af moralsk Natur, vil man see, at det Bør, som fremkommer i Logiken, kun er et
„saa bør det være", ei „saa bør \emph{vi} det;" og hiint Bør er derfor snarere et
Maa. --- Det egentlige Bør udtaler en Beherskelse som saadan, eller gaaer ud paa,
hvad der skal være Herskende hos den Beherskede, naar denne viser sig som saadan,
tilmed som attraaende og villende.\par}

{\small\emph{Anm.~4.} At Videnskaber af blandet Natur fremkomme, blive vi let vaer.
Vi have Exempler herpaa i Physik og Astronomie, forsaavidt i begge Empiriskt og
Mathematiskt gribe i hinanden. I Psychologien vil ved Siden af Det, som hører den
indvortes Erfarings Gebeet til, tillige Det, som hører den øvrige explicative
Ideephilosophie til, komme til Orde. I Æsthetiken maa under den ideephilosophiske
Iagttagelse og Udvikling bestandigen sees hen til de æsthetiske Grundbegreber, der
historiskt have gjort sig gjeldende. Moralphilosophie og Religionsphilosophie kunne
ikke gjennemføres, uden at empiriskt givne Forhold, hvortil jo endog hele
Syndigheden hører, komme i Betragtning. --- Men overalt kræver den grundige
Betragtning og Orientering, at man maa kunne blive sig bevidst, hvad der hører hvert
Gebeet til.\par}

% Rettelse (errata leaf, p.40 l.18, NB): print reads ``af det historiskt givne
% Sprog''; the NB erratum ``for det læs de'' gives ``af de historiskt givne Sprog''
% (plural: the historically given languages). Applied below.
{\small\emph{Anm.~5.} Vi saae her Psychologien at antage en blandet Natur ved at
komme ind paa det Ideephilosophiskes Gebeet. Men endnu i en anden Henseende viser
den sig som blandet, da udvortes Iagttagelse, nemlig af Andres Liv og Færd, altsaa
Noget, der hører til den udvortes Empiries Gebeet, her maa tages i Betragtning,
tilligemed Det, som den indvortes psychiske Iagttagelse, der bliver Grundlaget,
lærer os. --- Ogsaa Sprogphilosophien maa uddanne sig ved en Forening af empiriskt
Studium (af de historiskt givne Sprog) og psychisk Iagttagelse.\par}

{\small\emph{Anm.~6.} Endnu bliver her Historiens Philosophie at nævne som blandet,
da deri historisk Opfattelse og Forfølgelse maa gaae i Eet med en Iagttagelse, der
maa blive af psychisk Natur. Historiens Philosophie skal nemlig opklare
Menneskehedens hele Livsgang med Hensyn til, hvad der i Aandelighedens Henseende
deri er at bemærke, og den Udvikling, det i saa Henseende har taget.
Menneskehedslivets Grundmomenter og deres Bevægelsers Gang blive da at iagttage, men
saaledes bliver da Betragtningen psycho-biologisk.\par}

% Rettelse (errata leaf, p.41 l.14): print reads ``i en hvis Stat''; erratum
% ``for hvis læs vis'' gives ``i en vis Stat''. Applied below.
{\small\emph{Anm.~7.} En særdeles Opmærksomhed maae vi endnu henvende paa visse
Videnskaber, der danne egentlige Mellemled, idet de frembyde egentlige synthetiske
Eenheder af Noget, der paa udvortes historisk Maade er givet, og Noget, som hører
til den explicative Ideephilosophies Gebeet. --- Herved sigtes a) til den saakaldte
positive Aabenbaring og positive Theologie og b) til den positive Jurisprudents og
Politik. I den førstes Gebeet skal Noget, som væsentligen støtter sig paa historisk
Overleverelse, bringes til at leve og herske i os, og til at besjæle og opfylde os,
paa samme Maade, som den, paa hvilken Det, der hører Ideephilosophiens Gebeet til,
gjør sig baade ideebevægende og praktiskt gjeldende hos os, og hiint maa altsaa
exponeres paa en hertil svarende Maade. I Politiken, saavelsom i den i sig selv
dertil henhørende positive Jurisprudentses Gebeet, mødes, hvad idetmindste den
saakaldte Administration samt de til Lovgivningens Forbedring henhørende Bestræbelser
angaaer, bestandigen Hensynet til Det, hvad naturlig Ret og Statsomsorg fordrer, med
Hensynet til Det, hvad der i en vis Stat deels positivt er bleven fastsat og derved
kommen i Gang, deels historiskt har udviklet sig, og baade, hvad hint og hvad dette
i sin Conseqvents fordrer, maa paaagtes, og en Løsning af Opgaverne søges, der paa
eengang kan tilfredsstille i begge Henseender, altsaa baade i ideephilosophisk og i
empirisk-historisk Henseende.\par}

{\small\emph{Anm.~8.} At endeligen ogsaa explicativ og speculativ Betragtning kunne
forenes i een videnskabelig Udvikling og Fremstilling, maa endnu bemærkes, hvorpaa
da nærværende vor philosophiske Propædeutik kan være Exempel.\par}

{\small\emph{Anm.~9.} Af anden Natur end de hidtil omtalte Videnskaber af blandet
Natur, ere de mere almene Videnskaber, i hvilke Det, der viser sig som almeent til
Grund Liggende i flere forskjellige Gebeter, udsondres og betragtes for sig i sin
indre for disse Gebeter almeengyldige Grundnatur.\par}

{\small En saadan Videnskab er den almindelige Biologie, forsaavidt deri Livet i
Planter og Dyr, det indvortes sjælelige Liv i Mennesket, som Individ, og Livet i
Menneskeheden paa eengang og lige meget haves for Øie.\par}

\begin{center}\large\S~12.\end{center}
\markboth{\S~12}{\S~12}
\begin{center}Exposition af det givne Begreb om Philosophie: tre
Hovedsynspuncter herved.\end{center}

Til Indbegrebet af alt Det, vi have kaldet det Givne, og nu have stillet hen for os
i en Oversigt, træder Philosophien, --- hvorved vi nu, som hidtil, bestandig kun
mene den speculative, ikke den explicative, --- for at trænge hen til dets sidste
Grunde og Grundforhold og i Kraft af disses Erkjendelse at opklare det fra Grunden
af. Men idet den nu saaledes fremkaldes ved den altomfattende Randsagelsesaand, der
kommer i Bevægelse hos den granskende og grundende Subjectivitet, som, for at opnaae
en Philosophie, først maa philosophere sig til den, stiller Philosophien sig i
Henseende til Grundforholdet imellem det Sande og den Subjectivitet, hos hvilken det
Sande skal gjøre sig gjeldende, frem for os under tre Hovedsynspuncter, nemlig:
I.~med Hensyn til den erkjendende Subjectivitet og Philosophiens Betydning for
denne; II.~med Hensyn til det Subjectives og det Objectives Opgaaen til Eenhed;
III.~reent i og for sig, uden Hensyn til Det, som rører sig i den erkjendende
Subjectivitet.

\begin{center}\large\S~13, a.\end{center}
\markboth{\S~13a}{\S~13a}
\begin{center}Exposition af Begrebet Philosophie i første Hovedhenseende.\end{center}

Med Hensyn til I.~den erkjendende og til grundig Erkjendelse henstræbende
Subjectivitet lader den i Bevægelse værende Randsagelsesaand os allerførst see
Philosophien som a) den Grunderkjendelse, der skal gaae frem af en til Grunden
gaaende Kræven til Regnskab, gjennem hvilken Alt, hvert især, skal hjemle sig for os
i den Gyldighed, som det skal have for os, ved ligesom at opvise sin Adkomst til at
gjøre sig gjeldende. --- Men herved viser sig da b) Alt i sit hele Indbegreb som
Det, der frembyder et Problem, der skal løses, og Philosophien har derfor kunnet
fremstilles som den, der skal løse Tilværelsens Gaade, efter selv først at have
ladet en saadan Gaade opstaae og danne sig for os. --- Men idet hertil den
Bemærkning knytter sig, at de philosophiske Spørgsmaal ofte tage den Vending, at der
spørges om Mueligheden af Det, om hvis Virkelighed Ingen tvivler, ja hvis
Virkelighed forud maaskee formeligen oplyses, see vi Philosophien som den, der kan
siges c) at gaae ud paa at oplyse Virkelighedens Muelighed, hvilket da vil sige
Forklarlighed, Forstaaelighed, Begribelighed, hvor da Virkelighed og Tænkning sees
at mødes, som de, der søge Opgaaen til Eenhed og saakaldet Forsoning, hvilken da dog
tillige bliver en Udjevning og Forening af de mange imod hinanden trædende
Tilværelsens Momenter selv for den erkjendende Betragtning. --- Men da under al
saadan Tænkning Undersøgelsesaanden kommer i Bevægelse som Discussionsaand eller
Debatteringsaand, saa kommer d) Philosophien frem som den, der danner sig paa en
altomfattende Maade i Kraft af overvunden Dubitation, altsaa ifølge en Affirmation i
Kraft af overvunden Negation. (Cfr.\ om Erkj.\ og Gr.\ \S~20, P.\ 162.)

% Rettelse (errata leaf, p.43 l.24): print reads ``Jakobi''; erratum ``l. Jacobi''.
% Applied below (``som Jacobi kaldte det'').
Men saaledes vil da Philosophien fremstille sig for os paa den ene Side som den, der
skal føre Subjectiviteten til e) en fuldelig Kommen paa det Rene med sig selv
(Selbstverständigung, som Jacobi kaldte det), og paa den anden Side som den, der
skal føre til f) en fuldelig Kommen paa det Rene med Totaliteten i dens Totalitet og
derved til fuldelig Orienterethed i denne; hvoraf da g) en altomfattende
Tilfredsstillelse paa Erkjendelsens Gebeet maa blive en Følge.

{\small\emph{Anm.~1.} Kant begyndte sin Philosophies hele Værk (i Begyndelsen af
hans Kritik der reinen Vernunft) med det Spørgsmaal: hvorledes ere synthetiske
Sætninger \textit{a priori} muelige?, efterat han først havde oplyst saadannes
virkelige Existents og Realitet. Senere førte Dette til at spørge: hvorledes er
overhovedet en gyldig Erkjendelse muelig? men dette Spørgsmaal med det deri
liggende: hvorledes er en objectiv gyldig Sandsning muelig? havde allerede været i
fuld Bevægelse i den Cartesianske Skole, hos Malebranche og Leibnitz ligesom det
allerede i Oldtiden fuldt var traadt frem. --- Hvorledes er speculativ Philosophie
muelig? ville vi længere hen komme til at spørge. --- Hvorledes er menneskelig
Frihed muelig? hvorledes er det Onde mueligt? ere Spørgsmaal, som naturligen netop
antage denne Form.\par}

{\small\emph{Anm.~2.} Endnu kan det nominelle Begreb, som Fichte i den af ham i
Forening med Niethammer udgivne philosophiske Journal (Aarg.\ 1797, Hefte 1) giver
om Philosophie, samt det, han i sine seneste Forelæsninger i Berlin fremsatte, her
bringes i Erindring. Hist opstiller han Philosophien som den, der skal
forklarliggjøre hele Erfaringen overhovedet ved at deducere den \textit{a priori},
hvor man da, naar Philosophien fik fuldendt sit Værk, maatte faae at see, om Det,
den \textit{a priori} havde construeret, svarede saaledes til Erfaringen, at Opgaven
kunde kaldes løst. Da vi kun besidde hele Erfaringen i vor Viden, afveeg han herved
ikke fra sit tidligere Begreb om Philosophien, ifølge hvilket han havde givet
Philosophien Navn af Wissenschaftslehre, som den, der skulde gjøre den hele
menneskelige Viden i dens Heelhed, som Viden og Videns Indbegreb overhovedet, til
Gjenstand for Forklaring. I sine seneste Forelæsninger yttrede han, at enhver
Videnskab havde sit Phænomen --- han brugte just dette Udtryk --- at forklare. Det
Phænomen, Philosophien havde at forklare, var nu selve den menneskelige Viden
overhovedet eller i det Hele, als ein Wißthum. Dette Phænomen maatte nu først
exponeres, derefter forklares. Hiint gjorde han i Læren „von den Thatsachen des
Bewußtseyns", hvilken fremkom i den første Deel af hans Foredrag, --- Dette gjorde
han i den egentlige „Wissenschaftslehre". Man seer let, at han i Lighed med den
Maade, hvorpaa Kant begyndte sin Kritik, kunde begyndt med det Spørgsmaal:
Hvorledes er Viden overhovedet muelig?\par}

{\small\emph{Anm.~3.} Vi komme her til at mindes den ofte omtalte, men, som det
synes, i vor Tid almindeligen misforstaaede eller ikke kjendte Definition paa
Philosophie, som Wolf giver i sin \textit{Discursus præliminaris de philosophia in
genere} (foran for hans Logik). Den lyder i \S~29 saaledes: \textit{Philosophia est
scientia possibilium, qvatenus esse possunt.} Wolf siger om den i Anmærkningen, at
han først fandt den i Aaret 1703, at han i Aaret 1705 meddeelte den til den
breslauske Kirke- og Skole-Inspector Caspar Neumann, mod hvis Indvendinger han
forsvarede den i et privat Brev, men at han først bekjendtgjorde den i Fortalen til
sine \textit{Elementa Aërometriæ}. Denne giver imidlertid ikke den Oplysning, jeg
havde ventet at finde deri. Derimod oplyses Meningen af hiin Definition
tilstrækkelig i den anførte Discursus selv, vel ikke directe, men paa flere Maader
indirecte. De Ord, som det især kommer an paa, ere de sidste: \textit{qvatenus esse
possunt}. I en foregaaende Paragraph (\S~6) er \textit{cognitio philosophica}
defineret at være \textit{cognitio rationis eorum, qvæ sunt vel fiunt}, og man seer,
at han overhovedet betragter al rational Erkjendelse af Det, eller al Indsigt i Det,
der er eller skeer, som philosophisk, hvorimod den blotte Kundskab om det Factiske
opstilles som \textit{cognitio historica}, og den tredie af de tre Arter
Erkjendelse, han adskiller, den mathematiske, skal gaae ud paa den nøiagtige
Bestemmelse af de endelige Ting ved at bestemme deres Qvantitetsforhold. Naar han
nu, ved nærmere at omhandle denne Forskjel imellem historisk, philosophisk og
mathematisk Erkjendelse, blandt Andet siger i \S~17: \textit{Aliud est nosse factum,
aliud perspicere rationem facti, aliud denique determinare qvantitatem rerum}; naar
det \S~7 hedder: \textit{Differt cognitio philosophica ab historica, haec enim in
nuda facti notitia subsistit, illa vero ulterius progressa rationem facti palam
facit, ut intelligatur, cur istiusmodi qvid fieri possit}; og naar man nu i de
Exempler, han anfører, møder alleslags Phænomener, som Philosophien skal forklare,
saa seer man, at hans Mening maa være, at Philosophien skal forklarliggjøre
Virkeligheden, ved at vise, hvorledes Alt i den formaaer at være. --- Senere,
efterat der paa hiin Definition paa \textit{cognitio philosophica} er fulgt den
ovenfor anførte Definition paa Philosophien selv, hedder det i \S~31: \textit{In
philosophia reddende est ratio, cur possibilia actum conseqvi possint}, og strax
efter: \textit{In philosophia itaqve reddenda est ratio, qvomodo ea, qvæ fieri
possunt, actu fiant}; fremdeles hedder det \S~32: \textit{qvodsi ergo plura fuerint,
qvorum unum perinde possibile est ac cetera, philosophia docere debet, cur illud
potius fiat aut fieri debeat, qvam cetera}; ved hvilken Anledning Philosophien
omtales som \textit{rationes eorum, qvæ sunt vel fiunt, reddens}, og som den, der
derfor \textit{docere debet, qvænam sit ratio, ut in dato aliqvo casu unum potius
fiat, qvam alterum}. I \S~37 nævnes Philosophien som den, \textit{qvæ rationem
exponit, cur ea, qvæ sunt et fiunt, esse ac fieri possint}. I \S~46 endeligen finde
vi følgende Definition paa en Philosoph: \textit{Philosophus est, qvi rationem
reddere potest eorum, qvæ sunt, vel esse possunt}; og ved denne Leilighed nævnes
strax efter Philosophien som den, \textit{cujus est rationem reddere, cur possibilia
actum conseqvi qveant, et unum potius fiat in dato casu qvam alterum, qvod perinde
possibile erat}. --- Sammenholder man nu alt Dette og adskilligt Mere med hiin
Definition: \textit{Philosophia est scientia possibilium, qvatenus esse possunt},
saa seer man, at Vægten maa falde paa de sidste Ord, og vi see, at disses Mening maa
være at tilkjendegive hiint allerede Nævnte, at Philosophien skal oplyse, hvorledes
alt Det, som \textit{de facto} findes at være eller at skee i Naturen (hvortil ogsaa
efter \S~21 \textit{natura rationalis}, altsaa Bevidsthedslivet, henhører) formaaer
at være, eller hvilken dets reale Muelighedsgrund er. De virkelige Facta og
factiskt tilværende Ting har Philosophien ei med at gjøre, forsaavidt der er Tale
om, hvilke de ere; Dette hører den historiske Viden til; men oplyse Grundene til
dem, altsaa vise, hvad der fremkalder eller medfører deres Virkeliggjørelse, skal
Philosophien. Men Det, som medfører Virkeliggjørelse, er dog ikke Virkeligheden
selv, men er den reale Muelighedsgrund til samme. Saaledes maae vi uden Tvivl
udlægge hiin Definition. --- Philosophien bevæger sig i Tænkelighedens Gebeet --- der
tænkes frem og tilbage --- forsaavidt have vi med \textit{possibilia} at gjøre; men
nu gjelder det at bestemme, hvad der ud af al denne Tænkelighed hæver sig frem som
Det, der har Realitet, og udgjør det til Grund Liggende i Virkeligheden; hvilket da
ei er det concrete Virkelige selv, men er Det, hvori dette har sin Grund, og hvoraf
der altsaa er at indsee, hvorledes dette Virkelige formaaer at være (\textit{actum
conseqvi possit}). Det er altsaa Tænkelighedens Rige, der mødes med Virkelighedens;
hint skal oplyse dette, dette føres tilbage til hint og indoptages deri.
Virkeligheden er her uden Tvivl tænkt som den, der har et uudtømmeligt Muelighedernes
Rige til Grundlag: en Forestilling, som man i det 18de Aarhundrede ofte møder.\par}

\begin{center}\large\S~13, b.\end{center}
\markboth{\S~13b}{\S~13b}
\begin{center}Begrebets Exposition i anden Hovedhenseende.\end{center}

Men da nu dog de Fordringer, som her ud af Subjectivitetens Indre træde alt det
Givne imøde og kræve Tilfredsstillelse, ere af reen objectiv Natur, eller gjøre sig
gjeldende som de, der vel fremtræde i og gjennem Subjectiviteten, og det saa, at
denne ved deres Tilfredsstillelse skal komme til Klarhed og Beroligelse, men dog ere
det Algyldiges egne Fordringer, hvori dette taler i sin Algyldighed, saa kommer
Philosophien, II.~med Hensyn til denne det Objectives og det Subjectives Sammengaaen
til Eenhed, til at vise sig for os som den, der skal lade os see Alt a) i en
fuldkommen Samstemming med det sig hos Subjectiviteten selv i den hele
Erkjendelsesstræben rørende Algyldige, idet dette da træder frem som saadant, eller
--- med Fornuften, og det saaledes, at derfor, b) al Erkjenden, forsaavidt den er af
intelligent Natur, maa komme til at vise sig som den, der af Fornuften selv maa være
at deducere, saa at Philosophien altsaa ei kan Andet end sigte hen til en
altomfattende intelligent Erkjendelses Construction \textit{a priori}, og forsaavidt
viser sig som constitutiv i Kraft af fuld Fornuft. --- Men idet saaledes den
grundende og spørgende Betragtning kommer til at erkjende Alt i dets Harmonie med
Det, som constituerer denne Betragtens eget inderste Væsen, kommer den tillige til
c) en saa fuldelig Anerkjendelse af det Gyldige og Sande og derved tillige til en
saa fuldelig Assimilering eller Indoptagelse af Totaliteten af det i Tilværelsen
Givne, at dette derved d) taber sin Charakteer af Givethed, saa at vi blive hjemme
deri, og saaledes den Modsætning overvindes, som fandt Sted, saa længe det Givne
endnu blot stod for os som et Givet.

{\small\emph{Anm.~1.} Af ganske anden Art end den Tilegnelse og Indoptagelse, som vi
her omtale, er a) den, som skeer ved den Art af ideel Skuen, der lader os see det
Algyldiges og Alconstitutives indre Momenter træde ud og røre sig som det inderste
til Grund Liggende i den reelle Virkeligheds aldeles individuelle Skikkelser og
Livsbevægelser; hvilket deels skeer i en ideel Empiries deels i Poesiens og den
øvrige ideelle Konsts Gebeet, hvori da det til Grund Liggende viser sig i sin
magtfulde Nærværelse i det Concrete, paa ganske individuel Maade; fremdeles ogsaa b)
den, som paa sympathetisk Maade skeer ved vor egen interessefulde og villiefulde,
kjærlige og religiøse Tilknytning til og Leven i og for Tilværelsens Totalitet. Vi
have ogsaa i alt Dette Noget, hvorved eller hvorigjennem det Sande lige saa fuldt,
som gjennem Philosophering, kun paa en ganske anden Maade, kan sande sig for os; saa
at vel den Bemærkning, at det kan, ja maa det, bliver en philosophisk, men ei denne
Sandelse selv. Paa en saadan Sandelse er det, at Christus viser hen ved de Ord
(Joh.\ 7, 17): „Min Lærdom er ikke min, men hans, som mig udsendte. Dersom Nogen vil
gjøre hans Villie, han skal kjende, om Lærdommen er af Gud, eller jeg taler af mig
selv."\par}

{\small\emph{Anm.~2.} Vi have ovenfor bemærket, hvorledes al vor Erkjenden ei blot
gaaer ud fra, men ogsaa heelt igjennem væsentligen og nødvendigen hviler paa, et
\textit{Scimus, qvia accepimus} (see \S~3). Det er ogsaa af Vigtighed, at vi heelt
igjennem føle denne Givethed og føle det tilsvarende Fornemmelsesfulde i al sin
Styrke, og derunder føle Objectivitetens Magt (Herrens Magt) som saadan. Men nu
ligger dog heri strax --- og bliver uadskilleligt herfra --- et \textit{Scimus, qvia
facimus}, hvori det personlige Moment viser sig i sin Eenhed med, men tilmed ogsaa i
sin relative Modsætning til, hint universelle Moment. Men naar de, under denne
forenede Modtagen og selvegne Bestemmen, paa det første Stadium endnu i umiddelbar
Eenhed virkende og tilmed i Resultatet (den hele Antagen) fortabte Momenter komme
til at træde ud som saadanne, under den deels paa det Givnes Sammenhæng og Forhold
deels paa vor egen Antagens Gyldighed sig henvendende Reflexion, saa komme vi til en
Erkjenden, hvori vi befinde os, fordi vi selv have bekræftet os deri, og tilmed have
anerkjendt. Og nu kan det da hedde \textit{Scimus, qvia cognoscimus (nobiscum
noscimus)}. Herved viser sig da nærmest et \textit{Scimus, qvia intelligimus}. Men
nu have vi dog bemærket, at det ogsaa, baade i denne Henseende, og i Henseende til
det \textit{Scimus, qvia credimus}, der hører med ind under hiint \textit{Scimus,
qvia accepimus}, kan hedde: \textit{Scimus, qvia experti sumus; scimus, qvia coram
vidimus}, saa at der overalt kan vise sig en samtykket Erkjenden, ja at især den sig
selv bekræftende Tro er at kalde ei allene en samtykket men og en villet Erkjenden.
Men nu vil Forholdet, baade i denne Henseende, saa og igjennem Philosophie, kunne
gaae saavidt, at det kan hedde: \textit{Scimus, qvia sumus}; hvis det ellers tør
være tilladt at bruge det latinske \textit{sumus} paa samme Maade, som det
tilsvarende Udtryk bruges og den tilsvarende Tanke tænkes paa Dansk og Fransk, naar
Ordet construeres med et Sagsobject i Udtrykkene: det er mig, \textit{c'est
moi}.\par}

{\small\emph{Anm.~3.} Med Hensyn til, at Philosophie her er erkjendt at være
Videnskab i Kraft af fuld Fornuft, bringes i Erindring Hegels Definition i den
første Udgave af hans Encyklopædie \S~5: „Die Philosophie wird hiemit für die
Wissenschaft der Vernunft ausgegeben und zwar insofern die Vernunft ihrer selbst als
alles Seyns bewußt wird." Men herved er overseet, at al Væren væsentligen beroer paa
en Eenhed af Rationalt og Extrarationalt, Intellectuelt og Ikke-Intellectuelt, idet
al Væren paa den ene Side forudsætter den Magt eller Kraft, som realiserer eller
virkeliggjør det Begrebsmæssige, eller dog medfører, at det gjør sig gjeldende, paa
den anden Side forudsætter et sidste Substrat, hvori det gjøres gjeldende, som i et
Andet, end det selv.\par}

{\small Hvad der i vor Tid saa meget er bleven omtalt angaaende Tænkens og Værens
Eenhed, er der her Anledning og Opfordring til at tage i Overveielse. For det Første
synes herved, idet Sagen er taget subjectivt, at være meent, at vor Tænken, saasnart
den har fuld Gyldighed i sig, ogsaa har fuld objectiv Gyldighed (som da Astronomens
Calcul f.\ Ex., naar den er aldeles fuldkommen, ogsaa nødvendigviis maa blive
bekræftet af Virkeligheden); saa at altsaa vor Tænkning er at betragte som en
virkelig og gyldig Sands for den virkelige objective Intellectualitet. Forsaavidt
menes da den subjective og den objective Intelligents saaledes at være Eet, at hvo
der fuldeligen har den ene, \textit{eo ipso} har den anden. Dette er aldeles
grundet. --- Men nu er der dog dernæst ogsaa Tale om en Opgaaen til Eenhed af begge,
nemlig Tænken og Væren, i objectiv Henseende. En saadan kan nu kun være at forstaae
saaledes, at der menes, at paa ethvert Punct, hvor vi møde det Ene, er ogsaa det
Andet, men at dog ogsaa paa ethvert Punct foruden det Intellectuelle er tillige
noget Andet, end det, saa at Eenheden er en saakaldet synthetisk Eenhed af Tvende.
Eenheden er da ei af den Art, at man kunde sige: Hvo som har eller seer det
Intellectuelle, har og seer \textit{eo ipso} det Bærende; men man kan kun sige: Hvor
det Ene er, er ogsaa det Andet --- saaledes, at Eenheden ei er at kalde Identitet,
men snarere, om man vil, Formæling. --- Den, som fuldeligen har den subjective
Intelligents, lever \textit{eo ipso} i den objective Intelligentses Fylde; men Den,
som har den fuldelige Intelligents, eller Den, som har den fuldelige Fornuftfylde,
lever ikke \textit{eo ipso} i Synets, Hørelsens, Lugtens, Poesiens, Kjærlighedens,
Religiøsitetens Regioner o.\,s.\,v.; men man kan kun sige, at han lever tillige i
disse Regioner, forsaavidt nemlig Fornuftfylden ei fuldeligen kan haves, uden at
Livsfylden haves fuldeligen, og omvendt.\par}

{\small\emph{Anm.~4.} Angaaende den altomfattende Construction \textit{a priori} og
Philosophiens Constitutive bemærkes: Det maa fra først af ikke oversees, at her kun
er Tale om, at i Philosophien, naar den tænkes fuldendt, altsaa naar der sees hen
til Maalet, hvortil der sigtes, skal Totaliteten af den intelligente
Verdenserkjendelse constituere sig for os paa reen apriorisk Maade, men at det her
ei kan være Meningen, at Totaliteten af Verdensopfattelsen eller af det
Verdensbillede, vi bære i os, skulde nogensinde kunne constituere sig reent
\textit{a priori}. Naar vi overveie Dette, og overveie, at al grundig Indsigt maa
være en Erkjenden og Indseen fra Grunden af, men at den sidste Grund er det
Alconstitutive, der er det absolute \textit{prius}, og at dettes Manifestation i os
er Fornuften i os, saa vil Sagen vise sig klar og simpel. Vi kunne sige: Al Indsigt
overhovedet, altsaa og al Indsigt i Erfaringens hele Gebeet, maa baade komme af, og
vise sig som kommende af, Indsigtens Grundkilde, Fornuften.\par}

{\small Nærmere betragtende denne Sag, der har funden og finder saa megen
Indsigelse, men, naar blot al Misforstand i denne Henseende hæves, let vil sees kun
at være en ligefrem Følge af Philosophiens Begreb, bemærke vi Følgende: Philosophie
skal dog gaae ud paa at udgrunde og erkjende det Sande i dets Sandhed; er nu dette
allene at finde i Tingenes indre Væsen, som allene ved Forstand og Fornuft kan
erkjendes, saa maa Philosophie blive Fornufterkjendelse. Men Fornuften forholder sig
ikke saaledes til Det, den erkjender, som f.\ Ex.\ Synet til de belyste Gjenstande,
at den blot ligefrem skulde optage og fremstille det for sig, som det nu engang
findes at være. Fornuften kan Intet erkjende, uden ved at vedkjende sig det som
oprindeligen stemmende og harmonerende med den selv; hvilket følger af Naturen af
den Erkjenden, hvormed Indsigt og egentligt Begreb, altsaa og Prøvelse, skal være
forbunden. Men er nu hiint blot ved Fornuften paa denne Maade erkjendelige Tingenes
sande Væsen tillige det, der medfører, at disse Ting ere, som de ere, og udgjør det
altsaa disse Tings sidste Grund, og kan nu altsaa hverken det selv eller de
erkjendes i deres Sandhed, uden at det sees som deres Grund, og de som deri
grundede, det, som dem nødvendigviis medførende, og de som deraf følgende, saa
udfordres jo til den sande Erkjendelse, ikke blot at erkjende Tingenes Væsen,
hvilket er Gjenstand for den høieste Grundidee, men og at see, hvorledes alt Andet,
som deri er begrundet, er at udlede deraf, hvorved det først alt sees i sin sande
Betydning og i sit relative Væsen. Men at see, hvorledes det kan udledes af
Grundideen, er dog det Samme som, at vi selv i Forestillingen udlede det deraf, og,
forsaavidt Dette skeer paa en altomfattende Maade, deraf udlede, og ifølge deraf
opstille, et heelt System af Erkjendelser, hvori Alt sees --- men med det Samme ogsaa
sættes --- paa sin Plads. Men Dette er just at construere \textit{a priori}. Og er nu
Det, der saaledes at erkjendes i sit Væsen ved at erkjendes i sin Sammenhæng med det
høieste Sande, just Erfaringen og overhovedet alt det Givne, saa er det just denne
Erfarings hele Indhold, dette Givnes indre til Grund Liggende i dets hele Omfang,
der skal sees som Følge af Grundideen, altsaa udledes af den, altsaa og
gjenconstrueres \textit{a priori}.\par}

{\small Maae vi overhovedet indrømme, at Erkjendelsen er Det, gjennem hvilket det
Hele fuldeligen skal komme til Individets Bevidsthed, ved ligesom at gjenfødes i
denne som i et Afbillede, hvorved det Hele da kommer til at gjenleve i den Enkelte,
og maae vi nu endvidere erkjende, at det sande Første eller \textit{prius} i Tingene
er Det, der er den høieste Grundidees Gjenstand, det Evige, hvorledes kunne vi da
Andet, end erkjende, at dette paa samme Maade maa vise sig som det altbegrundende
Første i vor Erkjendelses Hele, som det er det i sig selv, dersom vor Erkjenden skal
gjenudtrykke hiint Tingenes Hele; saa at altsaa, ligesom ved hiint objectivt Første
alle Ting constitueres, saaledes ogsaa ved den Grundvirken, hvori dette Første
fremstiller sig i vor Bevidsthed, hiint indre Afbillede af alle Tingene maa paa
ethvert Punct constitueres. Men skulle vi nu heri og herved komme til at see og
erkjende Sandheden i dens Sandhed, saa maa ogsaa Afbilledet dannes i os under og
ifølge vor egen Indseen og ved Indseen ledede Bestemmen, altsaa maae vi selv i Kraft
af Grundideen udføre hiint det Heles Afbillede i vor Bevidsthed, men det er at
construere det i Kraft af Grundideen, der da maa vise sig som det altbestemmende
\textit{prius} i vort Erkjendelseshele. Ikkun forsaavidt Dette skeer, kunne vi sige,
at dette vort Erkjendelseshele virkelig og i Sandhed er overeensstemmende med hiint
Tingenes Hele, i hvilket jo Tingene ikke staae blotte og bare, hver for sig, men
alle kun ere til, forsaavidt de hænge sammen med det altbestemmende Første, og dette
er tilstede i dem.\par}

% Rettelse (errata leaf, p.54 l.20): print reads ``saameget''; erratum ``l. saa
% meget''. Applied below (``komme vi saa meget mere'').
{\small Høre vi imidlertid nu dog megen Opposition mod den Sætning, at Philosophien,
ifølge sit Begreb, gaaer ud paa en eftergjørende apriorisk Construction af det givne
Erfaringshele, ja læse vi endog hos Steffens: „Naturwissenschaft \textit{a priori}
ist der Tod aller Naturphilosophie", saa bringes vi naturligen til den Tanke, at her
ved Udtrykket \textit{a priori} er tænkt paa noget Andet, end vi derved tænke paa. Og
til denne Tanke komme vi saa meget mere, som vi see, at det er Naturphilosophien,
hvis Død al Naturvidenskab \textit{a priori} erklæres for at være. Vi kunne da enten
tænke os, at der ved hiin Naturvidenskab \textit{a priori} blot er tænkt paa en
saadan, der vil gaae ud fra visse Bevidsthedens saakaldte Kjendsgjerninger allene,
eller fra visse blot af deels logisk deels psychologisk Iagttagelse abstraherede
almindelige Forstandsbegreber og almindelige Sætninger om al Erkjendelsesvirksomheds
subjective Beskaffenhed, eller paa andre deslige Maader vil gaae ud fra noget
Subjectivt, som Mennesket bringer med sig, og vil gjøre gjeldende som nødvendige
Imaginationer, Noget, hvorpaa vi ikke blot hos Fichte finde eminente Exempler, men
som vi endog ikke kunne kjende Kants „Metaphysische Anfangsgründe" frie for. Eller
ogsaa kunne vi tænke os, at hiin Sætning, og andet Sligt, som af de fortrinligste
Philosopher høres, blot er Udtrykket af den høist grundede Erkjendelse, at vi, i vor
Tid, endnu ere alt for langt fra at være trængte saa vidt i Naturens Studium, at en
Bestræbelse efter, blot ifølge de Grundideer, hvortil vi nu allerede ere komne, at
ville construere sig et Naturens System \textit{a priori}, skulde kunne Andet end
være høist misligt.\par}

{\small Saa vist det er, at man i Moralen, ja endog i Psychologien, strax bør gaae
lige til den høieste Idee og bringe Alt i Forbindelse med den, --- thi den er det,
hvis Kraft skal besjæle Mennesket, og med hvilken han overalt skal søge at føle og
bringe sit Liv i Samklang, --- saa vist er det, at man derimod i Naturstudierne bør
holde sig nær til den levende Beskuelse af den givne Natur selv, og kun, ledet ved
den, efterhaanden berede Mueligheden af hiin Tilknytning til, og faste og bestemte
Nedstigen fra, og Deduction fra, den høieste Idee. Men en muelig Construction
\textit{a priori} bliver dog Maalet. Det er her ganske, som i Astronomien, hvori vi
have et talende Exempel paa, hvad Construction \textit{a priori} vil sige, og
hvorledes en grundig Tragten efter Erkjendelse kan og maa føre til en saadan. Hvad
vort Planetsystem angaaer, er Astronomien kommen saavidt, at den ud af sit
\textit{prius}, sin Grundidee, støttet paa visse Grunddata, kan \textit{a priori}
construere Totaliteten af alle Planetbevægelserne og af de herved indtrædende
Forhold. Det Samme tragter Lægen efter i sine Sphærer, og netop det Samme er det
ogsaa, Philosophien tragter efter i Henseende til Totaliteten af alle
Verdensforhold.\par}

{\small Men naar vi nu saaledes erkjende Dette at være den philosophiske Stræbens
høieste Formaal, saa opstaaer naturligen det Spørgsmaal, hvorledes vi nu da komme
til den altomfattende Grundidee, der skal være Grundlaget for den Construction
\textit{a priori}, hvorigjennem Alt, hvad der skal have Gyldighed for os, skal
constituere sig for os i sin Gyldighed. Og herved bliver da at bemærke, at denne
Grundidee, dette Grundlag, vistnok ei kan constituere sig for os paa den reent
aprioriske Maade, som vi f.\ Ex.\ finde i Mathematiken. Men herom, saavelsom om det
Spørgsmaal, om den tilsidst skal kunne indlyse for os i sin Gyldighed ved sig selv
allene, komme vi til at tale længere hen i \S~15 og 16.\par}

\begin{center}\large\S~14.\end{center}
\markboth{\S~14}{\S~14}
\begin{center}Begrebets Exposition i tredie Hovedhenseende.\end{center}

Endelig betragte vi Philosophien III.~reent objectivt, skuende hen paa dens Formaal
i og for sig, uden Hensyn til den erkjendende Subjectivitet og hvad der rører sig i
denne. Da see vi den at gaae ud paa at forfølge Alt saaledes, at det Alt maa vise
sig a) i sine sidste Grunde og Grundforhold, tilmed i sin sidste Grundethed; men
saaledes viser det sig da b) i sin Sandhed eller i sit sande Væsen, saa at der
altsaa sees, c) hvad Realitet og Betydning der egentligen er at tillægge det, altsaa
d) hvad det i sit inderste Væsen er. Men Dette medfører da, at e) alt Enkelt sees i
sin Indoptagethed i det Hele, saa at det Mangfoldige erkjendes i sin indre Eenhed,
og at det Hele træder frem for os i sin Heelhed; hvorved da den philosophiske
Grundstræben viser sig som en altomfattende Eenhedsstræben. Men under saadan Tragten
efter at see alt Enkelt indordnet i Tilværelsens Totalitet kommer da Philosophien
til at tragte efter f) en altomfattende Verdensanskuelse, hvori alle Ting kunne vise
sig for os i deres sidste Verdensforhold og i deres Verdensbetydning, og
Philosophien gaaer saaledes ud paa at vorde en speculativ Kosmologie.

Men denne Stræben maa nødvendigviis føre til at gaae videre. Thi den hele Verden
selv kan kun vise sig for os i sit rette Lys derved, at alt Endeligt sees g) i sin
Tilbageførelse til og Indoptagethed i det Alconstitutives, det Uendeliges og Eviges
Rige, hvorved da, omvendt, dette maa blive seet i sin Udtræden i det Endelige, som
det deri tilstedeværende Reale. --- Men Dette maa fremdeles føre os hen til h) det
Centrale i den hele Totalitet, og derved til Guddommen, og Philosophien maa da søge
sit sidste Grundlag i en speculativ Theologie, og den maa da blive en Eenhed af
speculativ Theologie og speculativ Kosmologie.

{\small\emph{Anm.~1.} Det naturligen her opstaaende Spørgsmaal, hvorledes dog en
Kosmologie, altsaa hele Philosophien selv, for os er muelig, ville vi længere hen
(\S~15, Anm.~2) komme til at besvare os.\par}

{\small\emph{Anm.~2.} Forsaavidt det Alconstitutive viser sig for os som det absolut
Begyndelsesløse og Tidløse, i hvis hele ideelle Fylde ingen Vorden, Omskiftning
eller Omdannelse endog kun er tænkelig, da al saadan netop allerede forudsætter det
Ideelle, og bestaaer i visse Realisationer af noget Ideelt, kalde vi det
Alconstitutive det Evige. Forsaavidt det derimod deels sees i sin indre
Uudtømmelighed som medførende en Utallighed af Determinationer, deels sees i sin
Relation til det Tiden Undergivne, som det heri Herskende og Levende, kalde vi det
det Uendelige. --- Endeligt derimod kalde vi Alt, hvad der, trædende frem, hvert som
Noget for sig, har sin bestemte Væren med sine bestemte Determinationer, saa at
derom kan siges: „Saa og saa er det; saa og saa er det derimod ikke", medens derimod
det Uendelige er alle Muelighedeers uudtømmelige Fylde. Men nu er da ethvert Enkelt,
og tilmed ethvert Endeligt, kun et Moment med i denne Fylde. Skue vi det retteligen,
saa skue vi deri det Uendeliges Fylde i eet af de Puncter eller Momenter, hvori det
er og rører sig, som i een af de til dets totale Organisation, eller til dets Rige,
henhørende Enkeltheder, hvilken imidlertid nu da sees i sin Fremgaaen af og i sin
Indoptagethed i dette Riges Totalitet, hvis Utallighed af Enkeltheder da, i deres
organiske Sammengaaen til Eenhed, netop er den Fylde, hvori det Uendelige er og
rører sig. --- Man vil heraf forstaae følgende Ord af Steffens i hans „Grundzüge der
philos.\ Naturwiss." Pag.\ 3: „Es giebt für das wahre Erkennen kein endliches Ding;
das Endliche wird vielmehr nothwendig in die ewige Form und sofort in das ewige
Wesen gesetzt, in welcher es aber mit diesem eins ist."\par}

{\small\emph{Anm.~3.} See vi nu tilbage, saa blive vi let vaer, hvorledes de her
anførte mange Momenter i Philosophiens Begreb og Sider at see den fra lære os, at
Philosophien vil kunne defineres ved Hjelp af flere forskjellige Begreber, som
gjøres til Grundlag, og at saaledes en heel \emph{Cyklus af Definitioner paa
Philosophie} vil kunne aabne sig for os, blandt hvilke nogle ligefrem ville blive at
henregne til de propædeutiske, andre derimod ville træde frem som saadanne, i hvilke
vor egen philosophiske Grundanskuelse lægger sig ind. Cfr.\ ovenfor \S~1, Anm.~1,
Pag.\ 3.\par}

\begin{center}\rule{0.18\linewidth}{0.4pt}\end{center}

%% ============================================================
%%  PHILOSOPHIENS GRUNDLAG OG DENS VÆRK.  (§§ 15--19, pp. 59--80; PDF 70--91)
%% ============================================================
\grouphead{Philosophiens Grundlag og dens Værk.}

\begin{center}\large\S~15.\end{center}
\markboth{\S~15}{\S~15}
\begin{center}Den philosophiske Grundidee.\end{center}

Vi skulle nu tale om, hvorledes Philosophien kommer i Stand. Men her opstaaer da
naturligen allerførst det Spørgsmaal, om der i at efterforske Tingenes sidste Grunde
og Grundforhold er at gaae nogen anden Vei, end en simpel fortsat og altomfattende
Explication og Forfølgelse af alt det Givne til alle Sider og i alle Retninger,
indtil man saaledes kommer til det sidste til Grund Liggende; eller om det med de
Tingenes første eller sidste Grunde, hvilke vi have nævnet, forholder sig paa en
anden Maade --- og om de ere at eftergranske ad en ganske anden Vei --- end den, der
finder Sted ved al det Slags Grunde og Grundforhold, som vi nu her, til Forskjel fra
hine, ville kalde de nærmere Grunde, hvilke i enhver Videnskabs Gebeet ere at
forfølge indtil deres eiendommelige første eller sidste Grundlag. Herpaa er at
svare, at, ligesom speculativ Philosophie er traadt frem og har gjort sig gjeldende
i Menneskeheden, som en Sag for sig, ved Siden af al anden videnskabelig Stræben,
saaledes har den ogsaa sit eget Grundlag og Udspring i Menneskeaanden, og det
forholder sig da med Spørgsmaalet om Tingenes sidste Grunde anderledes, end med
Efterforskningen af Det, vi her have til Modsætning kaldet de nærmere Grunde. ---
Under de mange givne Tings Betragtning træder der, udaf den tænkende Aands eget
Inderste (udaf Fornuften), Noget frem, som træder alt hiint Givne imøde, og nu
udgjør det aprioriske Grundlag for den hele speculative Philosophering, og udgjør
Det, der inderst rører sig i denne, og lader os see det hele Givne i det Lys, hvori
det viser sig for os, naar Spørgsmaalet om Tingenes egentlige Væsen og Realitet
kommer i Bevægelse. Og det er vel at mærke, at dette aprioriske Speculative, udaf
sit Grundlag og sin Grundbevægelse, ligesaavel træder det øvrige Aprioriske (det
Moralske, Religiøse med Mere) imøde, som det træder det i Sandselighedens Gebeet
Givne imøde. Dette aprioriske Speculative, der fremkalder og fremskynder den hele
speculative Philosophering, og som det immanente Grundbevægende rører sig deri paa
ethvert Punct, kalde vi den philosophiske Idee, den philosophiske Grundidee, der nu
skal gaae ind i alt det Givne, mediere det med sig, assimilere sig det og saaledes
% Rettelse (errata leaf, p.60 l.19, NB): print reads ``Enkeltheder og''; the NB
% erratum ``sæt Komma efter Enkeltheder'' inserts a comma (a reader's pencil slash
% marks the site in the scan). Applied below.
indoptage det i sig. Uden noget Saadant vilde overhovedet den hele Tænken ei tage
den speculative Vending, der jo forudsætter, at vi gaae ud over den hele Kjæde af
Enkeltheder, og da maa gaae frem af sin eiendommelige Impuls.

Vi see da heraf, at, uagtet der i den speculative Philosophie stræbes efter en reen
Fornufterkjenden, saa kan den dog ei danne sig uden under Ideebevægelsens bestandige
Conferents med alt det Givne i dets Givethed, ligesom det jo ogsaa bestandigen er
dette, der skal erkjendes i dets inderste Væsen, saa at Øiet derpaa bestandig maa
være henvendt, og den søgte Alopklaring væsentligen maa beroe paa en
Tilfredsstillelse af de fra denne Side kommende Fordringer.

{\small\emph{Anm.~1.} Det forholder sig med den speculative Idee, som med den
moralske Grundidee, og vi kunne opklare os hiin Idees Forhold til det Mangfoldige
ved at see hen til denne. Ved denne mene vi ikke noget vist bestemt, som i en
bestemt Sætning lod sig udtale, eller i noget vist bestemt Enkelt lod sig paavise,
men mene herved Det, som gjør, at vi i alle enkelte moralske Tilfælde dømme
saaledes, som vi dømme, og indsee, hvad vi indsee, og bevæges af Det, hvoraf vi
bevæges. Denne moralske Idee existerer altsaa kun i Totaliteten af alle moralske
Indsigter og Domme, som det paa ethvert Punct tilstedeværende Grundlag for det
moralske Erkjendelseshele, hvilket imidlertid kun kan komme til sin Indholdsfylde
under Ideens bestandige Conferents med Indbegrebet af den Deel af Menneskelivet og
Menneskeforholdene, hvormed Moralen har at gjøre. Den er nemlig fra først af kun den
endnu aldeles ubestemte og forsaavidt indholdsløse Idee om Det, som saaledes bør at
herske i os, at vi maae tragte efter at vorde beherskede og besjælede,
gjennemtrængte og opfyldte, deraf. Men naar der nu spørges, hvad det er, hvoraf vi
% Rettelse (errata leaf, p.61 l.15): print reads ``spørges hvad''; erratum ``sæt
% Komma efter spørges''. Applied above.
saaledes skulle vorde beherskede, og hvorledes i alle Henseender, saa vil det komme
an paa, at vi, næst efter at have undersøgt og forfulgt, hvad der ligger i det nu
Udtalte, dermed gaae ind i vor hele givne Leven.\par}

{\small Netop paa denne samme Maade har det sig nu ogsaa med den philosophiske
Grundidee, der fra først af ei er Andet, end den endnu ubestemte og forsaavidt
indholdsløse Idee om et alconstitutivt Alværende.\par}

{\small\emph{Anm.~2.} Vi ere her komne til det Punct, hvor vi kunne besvare os det
ovenfor (Anm.~1 til \S~14) fremkomne Spørgsmaal om Philosophiens Muelighed, med
Hensyn til, at den sigter til en altomfattende Verdensanskuelse, medens vi dog kun
kunne overskue saa ringe en Deel af den objectivt givne Verden. --- Det, som
constituerer Totaltilværelsen i dens Inderste, constituerer ogsaa vor Fornuft-Tilvær
i dens Inderste, og gjør sig nu i sin Totalitets Fylde gjeldende hos os indenfra
gjennem den philosophiske Grundidees Virken, og herved bliver speculativ Kosmologie
muelig.\par}

{\small\emph{Anm.~3.} Vi maae her see hen til det Sporadiske i Tilværelsen, hvilket
jeg andensteds har gjort opmærksom paa (først i Philos.\ Arch.\ og Repert., Pag.\ 256
ff., senere i min Psychologie, ny Udarbeidelse, \S~28 og flere Steder). Det er en
Følge af, at den organiserende Virken og Skaben i Tilværelsen --- som en
Totalorganisation betragtet --- paa ethvert Stadium i Udviklingsgangen træder ud i
alle sine Momenter paa eengang, fra utallige Substrat- og Subsistentspuncter af, for
at saaledes, idet alle Momenter gjøre sig gjeldende, hvert for sig, den hele
Mangfoldighed, der skal danne Organisationens Fylde, kan træde ud i sit Indholds
relative Adskilthed og Forskjellighed, for derpaa at gaae sammen til den Eenhed, der
kun som Eenheden af det saaledes Sondrede og relativt for sig Bestaaende bliver til
Fylde. Det Sporadiske er saaledes Organisationen paa Veien til at constituere sig,
paa hvilken dets Vei det hører den væsentligen til at gaae igjennem det Sporadiskes
Stadium, d.\ e.\ igjennem Tilværelsesmomenternes Debat. --- Tilværelsen danner sig
gjennem en Livs-Dialektik, lig den Dialektik, hvorigjennem (see længere hen) atter
Livet i vor Bevidsthed og tilsidst især Philosophien danner sig.\par}

{\small I Bevidsthedslivet kommer da nu saaledes den philosophiske Tendents til at
træde frem som en Tendents for sig i det hele System af Tendentser, og den er da
ikke fra først af, og kan ikke være, den i al vor Erkjenden inderst sig rørende, som
den ifølge sit Væsen tragter efter at være i videnskabelig Henseende, men maa nu
træde i Conferents med samtlige de øvrige. --- Verden oplader sig ei fra først af for
os allene fra sit inderste Grundlag af, for nu udaf dette, at bygge sig op i os ved
Magten af dette sit Centrale allene; den oplader sig tvertimod for os fra sin Fyldes
mange Puncter paa eengang, og ligesom nu den religiøse, den moralske, den poetiske
Tendents paa sin Maade, saaledes maa og den grundphilosophiske paa sin Maade først
tilvinde sig sin altomfattende Hersken. Det Centrale bliver saaledes selv kun et
Moment med i Totaliteten under dennes Henarbeiden til at constituere sig i os; ja det
er selv i sig sondret i et System af centrale Tendentser, der ogsaa indbyrdes have
at gaae ind i en Conferents med hinanden. [Hvad vi kalde Kræfter i Naturen er
overhovedet Verdenstendentserne, eller rettere det Alconstitutives Tendentser, i
deres Sondrethed, idet ogsaa \emph{De} danne en Organisation med Fylde, hvori hvert
Moment har sin relative Bestaaen for sig; --- en Uendelighed af Villier gaaer her
sammen eller skal gaae sammen til een fuldelig organiseret og saaledes paa
Mangfoldighedens Maade sig i sin Eenhed rørende Grundvillen og Grundhersken, i hvis
Rige imidlertid alle hine mere enkelte Grundtendentser vedblivende træde frem som
saadanne i deres Adskilthed; thi kun saaledes er Riget en Organisation med
Livsfylde.]\par}

\begin{center}\large\S~16.\end{center}
\markboth{\S~16}{\S~16}
\begin{center}Hvori den philosophiske Grundidee har sit Grundlag.\end{center}

Men nu opstaaer naturligen det Spørgsmaal, om da nu den philosophiske Grundidee selv
er at tage som noget Givet, hvorfra man har at gaae ud, eller om ogsaa den skal
godtgjøre sig for os i sin Gyldighed, ligesom alt Andet, og derved constituere sig
som Grundidee. Hertil maae vi svare, at den vistnok skal det, og at den altsaa selv
maa vorde underkastet den altomfattende Discussion og Prøvelse, som den selv sætter
i Bevægelse, idet den under denne Discussions Gjennemførelse selv maa affirmere sig
for os i sin Gyldighed. Men Dette kan den dog kun derved, at den viser sig som den,
der saa at sige, slaaer til, eller er i Stand til at præstere hvad den skal
præstere, ved virkelig at vise os Totaliteten af alt det Givne i et saadant Lys, at
det Alt derved bliver os klart og forstaaeligt i dets sidste Grunde, saa at Ideen
baade satisfacerer eller fyldestgjør sig selv for sig selv, saa og tillige
fyldestgjør os i vor Tilværelses Heelhed ved at fyldestgjøre os i Henseende alle vor
Tilværelses og alle den hele Verdenstilværelses Momenter.

{\small\emph{Anm.~1.} Men --- kan der siges --- naar nu dog selve denne hele saaledes
indtrædende Discussion og Debat kun kan føres under hiin Grundidees og tilmed under
en efterhaanden dermed given Grundanskuelses egen Indflydelse og Præsidium, kommer
da dog ikke denne Grundidee saaledes til at træde frem i Kraft af et oprindelig
Givet, og med Charakteren af et saadant? Hertil er at svare: Vistnok er det saa, og
derfor beroer al Philosophering paa en oprindelig philosophisk Sands og Impuls, fra
hvis Side der trædes alt det Øvrige imøde; men paa samme Maade faaer overhovedet
ethvert Moment i Liv og Tilværelse sin Stemme med under den altomfattende
Discussion, og Intet kan vises tilbage, der vil gjøre sig gjeldende som
medhenhørende Moment. Men nu kommer det an paa, at herunder nu Eenhedsideen saaledes
tilfredsstilles, at derved Det, i Kraft af hvilket denne Eenhed constitueres, eller
at den i denne Eenhed sig rørende Grunderkjendelse, som Fremtrædelsen af den deri
subsisterende og sig rørende Grundværen i og for vor Bevidsthed, virkelig kommer til
at vise sig som indoptagende alle disse Momenter i sig. Saaledes kommer da selve den
philosophiske Grundidee (selve Herskeren, eller Den, der skal arbeide sig frem til
at vorde det) til en Vexelkamp og Debat med den hele Mangfoldighed, eller med det
hele saakaldte Detail.\par}

{\small\emph{Anm.~2.} Vi føres ved det her Anførte tilbage til Forholdet imellem paa
den ene Side Philosophien, paa den anden Side --- deels Empirien samt dens forklarede
Gjenkomst i Poesiens og Konstens ideelle Gjenbilleder deraf --- deels Troen, altsaa
Religionen, samt hvad der videre fra den sympathetiske Erkjendens Side indstiller
sig. Det bliver forsaavidt Modsætningen imellem det Rationale og det Extrarationale,
eller rettere imellem den aldeles rationale og den paa en extrarational Grundvold
hvilende, forresten ogsaa i sin partielle Rationalitet fremtrædende, Erkjenden, vi
her komme tilbage til (cfr.\ \S~5, Anm.~5, P.\ 18).\par}

Vi bemærke nemlig, at, skjønt den speculative Philosophering ei blot maa see hen til
alt det herved Givne, men ogsaa maa agte paa det og respectere det og vide sig i
Samstemming dermed, skal den dog ei \emph{bygge} paa Erfaring eller paa Tro, ja den
kan det ikke engang, idet det kun er den selv, som stiller baade Erfaringens og
Troens Gebeter i det Lys for os, hvori de maatte sees, naar de skulde kunne danne et
saadant Grundlag.

Til nærmere Opklaring af denne Sag, hvad især Spørgsmaalet om Philosophiens Forhold
til Tro og Religion, og den Mening, at Philosophie tilsidst maa \emph{bygge} paa Tro,
angaaer, bemærkes her Følgende:

For selve Religionen og hvad det herved væsentligst kommer an paa, Religiøsiteten,
ligger Troen til Grund, ligesom der til Grund for Poesie ligger noget genialt
Poetiskt og til Grund for Musik ligger Gehør og den øvrige musikalske Sands samt
Skjønhedsfølelse og poetisk Følelse. Religionsbegreber endeligen kunne, ved
intellectuel Iagttagelse og explicativ Forfølgelse af det i og ved Religion og
Religiøsitet Givne i dets Sammenhæng og Conseqvents, udvikles og uddannes. En
explicativ Religionsphilosophie kan paa denne Maade komme i Stand, som den, der har
Troen til Grundvold.

Hvad derimod Forholdet imellem Tro og speculativ Philosophie angaaer, da skal paa
den ene Side Troen ikke saaledes søge sin Grundvold i speculativ-philosophisk
Overveielse, som om Philosophien var dens Basis, og dens hele Realitet altsaa
beroede paa de philosophiske Discussioners Udfald, men skal holde fast ved sin egen
heraf uafhængige Grundvold. Men paa den anden Side skal og kan nu dog heller ikke
Philosophien \emph{bygge} paa Tro, som om denne var dens egentlige Basis. Det hører
tvertimod til den egentlige Philosophies Værk, at den, ved Siden af meget Andet,
netop skal gjøre Spørgsmaalet om Troens Realitet, samt om Realiteten af Troens hele
Indhold, til Gjenstand for sin Overveielse og Discussion, og den kan altsaa, om den
end maa respectere Troen i dens Givethed, ligesom alt andet Givet, og maa lade den
faae sin indgribende Medvirkning og Stemme i den totale Discussion, ikke \emph{bygge}
paa Troen, da den maa \emph{bygge} paa Det, som skal stille baade Troen og alt det
Øvrige i det rette Lys. I Philosophien skal Troen med sit hele Indhold vise sig for
os i sin inderste og sidste Rationalitet eller Intelligenthed, eller i sin inderste
Ideesandhed; men Kilden til al saadan til det Sidste gaaende Rationalitet i vor
Erkjenden er Speculationen selv i dens Heelhed. [Jøvrigt forstaaer det sig af sig
selv, at vor Skuen af Troen i dens Ideesandhed ikke er selve den actuelle Tro;
forsaavidt kommer da Troen til at danne den blivende Modsætning til Philosophien, og
det paa collateral Maade; og allerede paa Erkjendelsens Gebeet kommer derfor det
Fuldelige i denne Henseende til at beroe paa en synthetisk Eenhed af Tro og
Philosophie. Paa lignende Maade forholder det sig med Philosophiens Forhold til det
Moralske, det Poetiske og Konstneriske, samt ogsaa med Philosophiens Forhold til
Empirien, i hvis Gebeet den fuldelige Erkjenden maa beroe paa en synthetisk Eenhed
af Philosophie og Empirie med lige Ret og Gyldighed fra begge Sider. Dog herom komme
vi længere hen, i \S~21, til at tale nærmere.]

\begin{center}\large\S~17.\end{center}
\markboth{\S~17}{\S~17}
\begin{center}Hvorledes Philosophie kommer i Stand. --- a) Explication i
speculativ Tendents.\end{center}

Spørge vi nu, hvorledes Philosophien kommer i Stand, saa vil Svaret blive, at den
naturligste Vei vil være den, at der A) gaaes ud fra Betragtninger over alt det
Givne, saaledes som det ad Explicationens Vei kan kjendes, men det saa, at Hoved- og
Grundpuncterne deri gjøres til Gjenstande for en speculativ Betragtning, der sigter
til at føre dem tilbage til een inderste altomfattende Grunderkjendelse, ved at agte
paa de Hentydninger, de maatte frembyde, hvilke imidlertid naturligviis kun kunne
fremstille sig som saadanne for selve den under saadan Betragtning sig fremhævende,
men endnu kun sporadiskt eller punctviis fremkommende, philosophiske Grundidee. Vi
kunne kalde Det, som saaledes danner sig, en Explication i speculativ Tendents.

Men nu finde vi under saadanne Overveielser Noget, hvilket vi overalt i alle
Videnskaber møde igjen som det fælleds Grundlag for dem alle. Dette er deels det
Logiske, deels det Ontologiske. --- Explicativ Logik og explicativ Ontologie maae
derfor danne den meest fundamentelle Deel af hiin Explication i speculativ Tendents,
da de gaae ud paa at opklare og forfølge de meest fundamentelle og tilmed meest
generelle Grundbegreber i den menneskelige Viden, hvilke vi ogsaa kunne betragte som
de første eller primitiveste Verdensforholds indre Indicier hos os. Men de dybeste
Grundforhold aabne sig først for os, naar vi i speculativ Tendents explicere os det
til Ideephilosophiens Gebeet Henhørende.

\begin{center}\large\S~18.\end{center}
\markboth{\S~18}{\S~18}
\begin{center}b) Speculation og Dialektik i væsentlig Eenhed.\end{center}

Til den egentlige speculative Philosophie skeer imidlertid dog først Overgangen,
naar B) under den altomfattende Conferents af Alt med Alt en altomfattende
Discussion og Debat af Alt imod Alt indtræder, hvorunder, idet alt Enkelt viser sig
for os i den relative Gyldighed, der bliver at tillægge det, og med den Betydning og
hertil svarende Plads, det kommer til at indtage i Totaliteten af vor Erkjenden, den
altomfattende Grundidee og tilmed Gudserkjendelse og Verdensanskuelse klarer sig
saaledes ud for os, at den nu kan blive Grundlaget for den søgte altomfattende
Opklaring --- Philosophien viser sig da her som den, der, saavel under sin Dannelse
som i sin endelige Udførelse, beroer paa en \emph{væsentlig Eenhed af Speculation og
Dialektik}.

Ved Speculation forstaae vi den dybere Grunden og aandelige Skuen, hvorved den sig
fremhævende philosophiske Grundidee bestandigen holder sig nærværende for Sjælens
Øie, saa at den bestandigen og paa ethvert Punct gjør sig gjeldende, og herunder
udvikler og uddanner sig, saaledes som det ovenfor (\S~15, Anm.~4, Pag.\ 60--61) ved
Exemplet af den moralske Grundidee er bleven oplyst. Speculationen svarer saaledes
til Det, vi ellers kalde Videnskabens Aand. --- Ved Dialektik derimod forstaae vi den
philosophiske Tankeforfølgelse, hvorved, under Indflydelsen og Herredømmet af hiin
speculative Grunderkjendelse, enhver Gjenstand, ethvert Punct og Moment, discuteres
og debatteres saa, at det kan vise sig, hvad hver Gjenstand, hvert Punct og Moment,
skal gjelde for at være, hvad Betydenhed og Væsenhed, hvad Realitet og Ikke-Realitet,
der bliver at tillægge samme. Dialektiken kan her siges at være den sig i Forhold til
det Mangfoldige sættende og dette bestemmende og efter sig selv formende Speculation
selv; (saaledes har jeg ogsaa kaldet Dialektiken i mit Skr.\ om Erkj.\ og Gr.\ \S~14,
hvortil, saavelsom til sammes \S~13, her i det Hele henvises). Alt Enkelt kommer
herved, baade i Virkelighedens og i Begrebernes Sphære, til at træde frem i sin
Relativitet, visende ud fra sig hen til Andet, hvorpaa det beroer, hvoraf det
afhænger, og hvori det kun er et Moment, bestandigen tillige selv trædende frem med
forskjellig Gyldighed, altsaa ogsaa Ikkegyldighed, efter Relationerne, hvori det
kommer til at skulle eller ville gjøre sig gjeldende.

{\small\emph{Anm.~1.} Jeg henviser endnu, hvad Speculationen især angaaer, til min
Afhandling om den intellectuelle Anskuelse i Philos.\ Arch.\ og Repert.\ P.\ 125 ff.;
angaaende Dialektiken see sammesteds P.\ 387 ff., samt min Logik, \S~90.\par}

{\small\emph{Anm.~2.} Her fortjene følgende Ord af Steffens i hans „Grundzüge der
philos.\ Naturwiss.", Pag.\ 36, at erindres: „Wie das Wissen speculativ wird durch
die Anschauung der Universalität, so wird die Speculation wissenschaftlich durch die
Anschauung der Individualität."\par}

{\small\emph{Anm.~3.} Det Speculative reent allene, adskilt fra det Dialektiske, ja
fra det Rationale overhovedet, selv i dettes explicative Form, have vi paa en egen
Maade i det Mystiske (see den nysnævnte Afhandling, Pag.\ 235). --- Det dialektiske
Moment allene, uden speculativ Gehalt, have vi, reent og bart, i Sophistiken og
Scholastiken, med en bestemt Tendents derimod, og saaledes med en vis Gehalt, kun
ingen videnskabelig, men blot en praktisk, have vi det i de Gamles Skepsis. --- Ved
Parmenides kunde man maaskee ogsaa have Anledning til at spørge, om ei hos ham det
Speculative træder frem uden det Dialektiske.\par}

{\small\emph{Anm.~4.} Endnu er at bemærke, at philosophisk Gehalt, i rational
Fremtræden, som dybtskuende Erkjendelse i Henseende til de inderste Verdensforhold i
Aandelighedens Henseende, kan træde frem uden videnskabelig Form i Poesien, i
anskuelsesfulde Yttringer af den speculative Grundidee; thi ogsaa denne kan paa en
vis umiddelbar genial Maade bryde igjennem i saadanne, ifølge den vakte Sands for
samme, uden at have gaaet Tankeforfølgelsens og Discussionens Vei, om end ikke uden
et vist Øie for Livets Dialektik. Allerede Platon bemærkede disse og andre det
høiere sig fremarbeidende Livs geniale Virkninger som dem, der førte til en vis
rational Erkjendelse, uden at have dannet sig paa Rationalitetens egentlige Maade. En
vis oprindelig Sympathie imellem Makrokosmen og Mikrokosmen og en herfra
udspringende sympathetisk Bevægethed af visse af Ideefyldens Frembrud, ifølge
Fornuftens umiddelbare Blik, kommer her tilsyne i den intelligente Erkjendens Gebeet,
altsaa udenfor det Gebeet, jeg har kaldet (cfr.\ ovenfor \S~5, Anm.~4) den
sympathetiske eller hengivelsesfulde Erkjendelses Gebeet.\par}

{\small\emph{Anm.~5.} Gjennem den altomfattende Discussion, saavel i positiv Retning,
forsaavidt den bestemmer hvert Begreb sin relative Betydenhed og tilmed sin Væren og
saa at sige sin Rolle i det hele System, som og i negativ Retning, forsaavidt
Relativiteten, altsaa og Overgangene, heri sees og forfølges, viser Dialektiken sig
som den, hvori Stringents og Fluiditet træde frem i synthetisk Forening. Hvert
Begreb, idet det holder sig fast i sin eiendommelige Betydning og i sammes
Conseqvents, holdes dog tillige bestandig aabent for en vis Bestemmelighed, saa at
det bestandigen kan vende sig forskjelligen, og under sin Anvendelse bestandigen,
medens Andet bestemmes i Kraft deraf, selv vorde underkastet deels Prøvelse deels
Modification, idet Intet tør opstilles som fast forud for det Andet, men kun alt
Enkelts Gjennemgaaelse og Bestemmelse giver hvert Enkelt sin Fasthed og Betydenhed.
--- En modsat, den mathematiske, især den geometriske, mere lignende Maade at tage
Philosophien paa har man kaldet Dogmatismus. Mod en saadan opstillede Kant sin
Kriticismus, fordrende en altomfattende Prøvelse af Erkjendelsesvirksomheden selv
forud. --- Dogmatikere modsættes ellers Skeptikere. Men Dogmatismen finder sin
egentlige Modsætning i Dialektiken.\par}

{\small\emph{Anm.~6.} En Livets Dialektik er ovenfor (i Anm.~4) bleven nævnet.
Herved menes Livsmomenternes Debat, under den Livets Kamp, hvorigjennem det Gyldige
efterhaanden kommer til at vise sig i sin Gyldighed, det Relative i sin Relativitet,
det Ureale i sin Urealitet, og hvorved saaledes tillige for os en Livets Skole danner
sig, idet det efterhaanden for den Tænksomme og Sindige klarer sig ud, hvor det ægte
Tilhold i alle Henseender er at finde, hvor ikke. En saadan Livets Dialektik kan da
tillige vise sig som en Livets Dikaiosyne, en Livets Nemesis.\par}

\begin{center}\large\S~19.\end{center}
\markboth{\S~19}{\S~19}
\begin{center}Philosophiens Udtræden i c) et philosophiskt System.\end{center}

Ved den speculative Dialektiks Gjennemførelse, under hvilken Philosophien viser sig
som den, der lige meget beroer paa Gehalt og paa Form, banes Veien til C) at
uddanne den philosophiske Erkjenden til en systematisk Erkjenden, der saa derpaa kan
gaae til at fremstille sig i et philosophiskt System, men maa finde sit Complement i
sin Modsætning, det frie Judiciøse, og maa blive af organisk Natur, idet deri heelt
igjennem alt Enkelt lige meget maa bekræftes ved og vorde bekræftende for det Øvrige,
saa at Erkjendelsen af hvert Enkelt hviler i Totalanskuelsen, medens denne dog atter
maa hvile i alt Enkelts Opgaaen deri.

{\small\emph{Anm.~1.} Af saadan organisk Natur er ogsaa al explicativ Videnskab i
sit Gebeet, forsaavidt den nemlig ikke vil blive staaende ved at være en blot ordnet
Kundskab om det factiskt Givne (under hvilket da ogsaa Naturlovene, forsaavidt de
blot ere udfundne, ei forklarede, blive at henregne), men vil gaae til en vis
Grunderkjendelse, der skal give Indsigt i det saaledes Givne. Den blotte Explication
er derimod ei af saadan organisk Natur; ei heller er Mathematiken det, da det Enkelte
deri kan staae fuldkomment fast, inden det Øvrige er opgaaet for Erkjendelsen, saa at
det Enkeltes Erkjendelse ikke hviler i Totalanskuelsen. Men derimod kan nok noget
Dialektiskt komme tilsyne i Mathematiken, naar Bestemmelsernes gjensidige Afhængighed
af hinanden viser sig, og hvert Moment sees at spille sin relative Rolle i en
Totalfunction; og paa denne Maade kan da ogsaa indenfor en vis Hovedfunctions Sphære
noget Organiskt vise sig. Vel vil Schmidten i sit interessante Skrivt „om
Mathematikens Væsen" Pag.\ 31, at Mathematiken heelt skal have hiint Organiske; men
jeg finder ikke, at han oplyser denne sin Paastand, endsige retfærdiggjør den.\par}

{\small\emph{Anm.~2.} Af Philosophiens organiske Natur indsees det nu ogsaa, hvorfor
den ei kan træde frem med den Stringents og absolute Vished, som Mathematiken,
hvilket da hænger sammen med Philosophiens høiere Charakteer. Og her vil det da
tillige være at see, hvorledes de mangfoldige Stridigheder mellem Philosopherne,
angaaende endog det allerførste Grundlag, ere at betragte: Stridigheder, som man
forresten paa Religionens Gebeet har havt i lige saa stor Mængde. Den altomfattende
Debat af Alt mod Alt, der skal foregaae som den, hvorigjennem Grundanskuelsen først
skal klare sig ud for os, og som derfor hos ethvert Individ, hos hvilket
Philosophien fra Grunden af skal danne sig, maa foregaae og gjennemkæmpes paa Ny,
møde vi i det Større i hine Stridigheder, og see saaledes mange store Tænkere blandt
Philosopherne som Repræsentanter for eet eller andet Hovedmoment, hvilket nu træder
frem gjennem dem med eensidig Styrke og i eensidig Retning. Intetsteds sees det mere,
end i de philosophiske Debatters og i de religiøse Bevægelsers Sphære, at vor
aandelige Leven i det Hele endnu befinder sig i sin første Dannelses- og
Constitutionsproces, og derfor endnu først søger hiint faste, uanfægtelige Grundlag
og hiin sikkre Vei og Retning i sin Gang, som Kant, i Fortalen til sin Kritik der
reinen Vernunft, i hvilket Værk han haabede at etablere et saadant Grundlag og en
saadan Vei for Philosophien, bemærker, at Mathematiken allerede havde opnaaet, da
Thales, eller hvo det ellers monne have været, udfandt Loven for den ligebenede
Triangel, ligesom Kant ogsaa mener, at den formelle Logik allerede før Aristoteles
maa antages at have opnaaet sin Fasthed, og at den empiriske Naturlære, især ved at
gaae Experimentets Vei, er kommen til at naae en fast Videnskabs sikkre Gang. --- Det
maa naturligviis være anderledes, hvor den enkelte Erkjendelses fulde Gyldighed
væsentligen beroer paa, at det altomfattende Totale staaer fast for vor Anskuen i
Henseende til samtlige sine Hovedpuncter eller Hovedmomenter, og hvor altsaa Regelen:
\textit{Totum est parte sua prius}, ogsaa bliver en Regel for Indsigten. --- Mange
have troet, at Fasthed i Henseende til den Betydning, der tillægges Ordene, maatte
bidrage meget til Enigelse; men det er egentlig en Fasthed i Henseende til
Begreberne, det kommer an paa; men bag enhver Begrebsbestemmelse, naar den ikke skal
være blot nominel og altsaa lade Sagen selv staae uopklaret, ligger i
Philosopheringens Gebeet hos den dybere Tænkende den hele Grundanskuelse af
Verdensforholdene, ligger endog en Grundtendents, som en Grundvillen. --- Endnu kan
her bemærkes, at de philosophiske Stridigheder have givet en naturlig Anledning til,
at ved Siden af en vis Skepticismus en vis Eklekticismus har kunnet danne sig.\par}

{\small\emph{Anm.~3.} Angaaende Forskjellen imellem systematisk Erkjenden og sammes
Fremstilling i en bestemt systematisk Form, samt Forholdet imellem det saaledes
Systematiserede og det frie Judiciøse, see: Om Erkj.\ og Gr.\ Anmærkningerne til
\S~15, Pag.\ 135, ff.\par}

{\small\emph{Anm.~4.} Hvad Philosophiens saakaldte Inddeling angaaer, da kan herom
Intet paa dette Sted directe læres, da den kommer til at beroe paa den hele
philosophiske Ideegang, saaledes som denne, som et Slags Tilværelsens egen
Manifestationsgang, vil danne sig for den Philosopherende under Udførelsen af selve
Philosophiens Værk. Derfor skal her kun i en Anmærkning være Tale derom.\par}

{\small En Overveielse og Forfølgelse af Philosophiens subjective Hovedopgave vil
man, saaledes som Kant i Grunden gjorde det ved sin Kritik, og som man hyppigen har
gjort det ved at indlede sin Philosophie ved en saakaldet „Erkenntnißlehre", meget
vel kunne lade træde op som første Led, skjønt vi derved allerede føres midt ind i
hele Tilværelsens Midte, saa at Philosophiens Cykliske, grundet i dens Organiske,
herved strax fuldeligen kommer tilsyne. Vi gjøre derved fra først af Rede for os,
forsaavidt vi betragte den subjective Hovedopgave som indoptaget i den objective; men
vi maae da vistnok siden paa sit Sted lade Det, vi saaledes indledningsviis have
taget for, indordne sig i Totaliteten af den Tilværelsens Manifestationsgang, vi
siden forfølge.\par}

{\small Dernæst maatte, som mere egentlig reel Indledning i Philosophien (forskjellig
fra dens Propædeutik), en Gjennemførelse af den forhen (\S~17) nævnte Explication i
speculativ Tendents indtræde, for at den Grundanskuelse, der siden skulde danne
Grundlaget, kunde, gjennem den hermed forbundne Conferents af Alt med Alt og Debat af
Alt mod Alt, altsaa gjennem Dialektik, constituere sig forsaavidt for os, at den nu
derpaa kunde lægges til Grund, og tillige fremdeles, under sin Udførelse og
Gjennemførelse, staae sin Prøve. Saaledes fik vi da en Gjennemførelse af den
Philosophering, hvorved man skal philosophere sig til den sidste altomfattende
Grunderkjendelse, og erholdt herigjennem allerede Philosophien selv paa indirect
Maade.\par}

% Rettelse (errata leaf, p.75 l.15): print reads ``Tilværelsesphære''; erratum
% ``læs: Tilværelsessphære''. Applied below.
{\small Selve Philosophien maatte da vel begynde med en Fundamentalphilosophie, der
gav en formelig Fremstilling af det hele Grundlag; derefter vilde da Tilværelsen være
at gjennemgaae saaledes, at Grundanskuelsen gjennemførtes gjennem de samme
Tilværelsens Sphærer, som man i hiin reelle Indledning foreløbig havde betragtet,
hvorved da først Læren om de almene Tilværelsens Grundforhold, --- en saakaldet
Ontologie, --- dernæst en Naturphilosophie (forsaavidt man var i Stand dertil) og
derpaa en Aandsphilosophie maatte fremkomme. I dennes Gebeet maatte man fra
Naturgrundlaget for det sig i Bevidstheden opladende objective Aandelige gaae over
til at betragte dette Aandelige selv. Og her vilde da vel først Pladsen være for den
egentlige Logik, forsaavidt denne dog er Læren om det Logiske i dets Fremtræden i
erkjendende Væsener, samt for Philosophiens Philosophie, der da vilde stille sig paa
lige Linie med Læren om det Logiske. Hertil vilde da komme den speculative
Betragtning af det Æsthetiskt-Poetiske, det Moralske og det Religiøse. Ved de to
sidste vilde vi derpaa føres til at fremstille den herved styrede og ordnede
Menneskeleven, baade i individuel Henseende og i Henseende til de Enkeltes Sammenvært
til at danne de større organiske Heelheder, Stat og Kirke, hiin som gjeldende for det
jordiske Standpunct, denne som visende ud over den Tilværelsessphære, hvori vi her
paa Jorden befinde os, saa at her Læren om Udødeligheden vilde finde sin Plads.\par}

{\small\emph{Anm.~5.} Philosophiens saakaldte Methode i systematisk Henseende, eller
Ideegangens hele Art og Maade, vil naturligviis saaledes beroe paa, hvad der under
Udførelsen af selve Philosophiens Værk vil vise sig, at ogsaa herom Intet forud kan
læres, Andet, end at man i ethvert Gebeet maa lade den gaae den Gang, Sagen selv hver
Gang vil fordre, uden heelt igjennem at lægge an paa at gjennemføre en vis
Grundtypus.\par}

{\small\emph{Anm.~6.} Det er især ved at træde frem som systematisk, at den
speculative Philosophie viser sig som constitutiv. Cfr.\ ovenfor \S~13 b, Anm.~4,
Pag.\ 51, ff. Til det hist Bemærkede bliver her Følgende at tilføie:\par}

{\small Den philosophiske Grundidee og den deraf fremgaaende Grundanskuelse skal
uddanne sig for os og bekræfte sig for os som den i Henseende til Totaliteten af det
Givne gyldige; og saaledes kan den da hverken fra først af danne og constituere sig,
ei heller lede til en heel Philosophies Gjennemførelse og Constitution, uden gjennem
og formedelst den altomfattende Conferents med alt det Givne, eller med alt Det, hvad
den, ifølge sit Begreb, skal tjene til fra Grunden af at opklare for os. Alt det
Givne skal jo indoptages saaledes i Grundanskuelsen, at det alt derved kan sees i sin
inderste Rationalitet eller Fornuftighed; men saa maa det da og paa ethvert Punct
vise sig, at det nu virkelig er det os i Verden Givne, --- Det i hvis Midte vi ere,
leve og røre os, --- hvilket nu saaledes sees i sin Rationalitet, sees i Grundideens
Lys, og sees som indoptaget i Ideens Rige; men saaledes maa da Philosophien, ei blot
under Henarbeidelsen til sit Grundlag, og medens den søger dette og tilmed søger
Udgangspunctet for sin Deduction og Constitution, men ogsaa i sin hele Udførelse og
Gjennemførelse i og gjennem den systematiske Heelhed, paa ethvert Punct være sig sin
Samstemming med Virkeligheden vis, som den, hvorpaa dens egen reale Gyldighed beroer.
Ligesom den kun kan danne sig under en altomfattende Conferents af Alt med Alt og
Debat af Alt mod Alt, saaledes kan den ogsaa, under Udførelsen og Gjennemførelsen af
den hele systematiske Erkjenden, kun træde frem og udfolde sig i sin altomfattende
Gyldighed ved at vise sig som den, der har optaget i sig den hele Conferentses og
Debats Resultater, som sine Momenter. Ligesom den kun vorder den, den vorder, ved
dette sit saaledes tilvundne Indhold, saa maa den og fremdeles bære dette i sig, som
det for den givne Virkelighed gyldige Indhold eller gyldige Indbegreb af Momenter;
men saa maa den da ogsaa i Systemet paa ethvert Punct træde frem med denne sin
Gyldighed, altsaa i og med sin Samstemming med Virkelighedens Rige. Ellers vilde jo
ei heller den Udsoning og Opgaaen til Eenhed indtræde, ifølge hvilken vi skulde
bevæge os i Ideens Rige, som i Virkelighedens, og i Virkelighedens Rige, som i
Ideens. (Nemlig hvad Virkelighedens Rationalitet angaaer; thi kun om denne er her
Tale. Om det Extrarationale, der fra et Rationalitetens Rige gjør Ideeriget til et
Virkelighedens Rige, tales her ikke; her handles kun om Eenheden af Ideeriget og
Virkelighedens Rationalitet).\par}

{\small Og her kunne da tvende Spørgsmaal besvare sig for os. Det ene er det ---
underlig nok --- i vor Tid fremkomne Spørgsmaal: i hvad Forhold dog Philosophien
staaer til Liv og Virkelighed. Det andet er det, som betreffer den i vor Tid
fremkomne Paastand, at Philosophien skal begynde aldeles forudsætningsløs.\par}

{\small Hvad denne Paastand især angaaer, da maa man tvertimod opstille den Sætning:
at Philosophien, idet den begynder at danne sig, maa have Totaliteten af alle
Forudsætninger for sig, og ikke allene maa uddanne sig under det herved fremkommende
Vexelvirkningsforhold, men ogsaa, naar den nu skal constituere sig i det
philosophiske System og begynde sin Deduction og Construction \textit{a priori}, maa
træde frem med det under hiint Vexelvirkningsforhold dannede og uddannede Grundlag,
med den gjennem dette Vexelvirkningsforhold til sin Indholdsfylde komne Grundidee,
der har indoptaget og opløst alle hine Forudsætninger i sig, og nu viser sig som den,
der bærer en hertil svarende Fylde i sig, altsaa paa ethvert Punct træder frem i sin
Samstemming med Virkeligheden. (Iøvrigt henviser jeg her til mit Skrivt „Bemærkn.\ og
Undersøg.\ betreffende Hegels Philosophie", der er et særskilt Aftryk af nogle i
Maanedskr.\ f.\ Literatur 10de Aargang (1838) indrykkede Artikler).\par}

{\small Selv om man gik op til det høieste Punct og tænkte sig Philosophien at være
den guddommelige Albestemmen lig, der jo maatte synes at være uden al Forudsætning,
vilde man ei komme til andet Resultat. Thi for det Første blev at bemærke, at den
guddommelige Albestemmen og Alconstitueren ingenlunde kan tænkes som den, der, naar
den skulde begynde i Tiden, maatte begynde forudsætningsløs; thi forud for al saadan
Begynden maatte jo den hele guddommelige Fylde gaae som det absolut Begyndelsesløse, i
hvilket alt Enkelt er og hviler i Totaliteten. Tænktes da Philosophiens Begyndelse i
Lighed hermed, saa maatte jo ogsaa dens Fremtrædelse allerede ved det første Skridt
have Totalideen med dens hele Indhold til sin Forudsætning. Og vil man nærmere gaae
ind i det guddommelige Væsen, og nu heri især see hen til Viisdommen, som den, hvis
Gjenbillede Philosophien især skulde være, og spørge sig, om nu ikke dog den
guddommelige Viisdom udaf sig selv allene constituerer den totale Tingenes Orden, og
om ei paa en lignende Maade den til Fuldelighed komne Fornuft maatte reent udaf sig
selv allene constituere Philosophien, saa maa Svaret blive, at den guddommelige
Viisdom ei kan tænkes forudsætningsløs; thi den har den hele Guddomsfylde til sin
Forudsætning, og nærmere maa især Kjærligheden vise sig som dens Forudsætning, eller
som Det, vi strax maae tænke i Eet med den. Men paa lignende Maade maa Philosophien
strax fra først af gaae i Eet med reel Livsopfattelse og reel Livsanskuelse, og have
den hele fuldelige Leven til sin Forudsætning.\par}

{\small Exempel paa en Videnskab, som fremfor al anden kan siges i en vis relativ
Forstand at begynde forudsætningsløs, have vi i den rene Mathematik, der begynder med
visse første Observationer af visse Mueligheder, og nu ifølge heraf selv constituerer
sig Det, den videre forfølger. Hvis man vil indvende, at den dog forudsætter den
almindelige menneskelige Tænkning med de for samme til Grund liggende almindelige
Grundbestemmelser eller saakaldte Kategorier, og almindelige Grundfordringer, der uden
videre komme til at gjelde og gjøre sig gjeldende i den mathematiske Tænkning, som
Tænkning, medens derimod Philosophien gjør selve disse allerførste Forudsætninger til
Gjenstande for en Undersøgelse og Kritik, --- saa er hertil at svare, at selve den
Videnskab, der skal undersøge disse første Forudsætninger for al Tænkning, dog selv
ikkun kan komme i Stand formedelst Tænkning, og altsaa netop lige saa fuldt, som den
mathematiske Tænkning, maa lade de i al Tænken, blot som Tænken, liggende
Grundforudsætninger gjelde og gjøre sig gjeldende just i den Tænkning, ved hvilken de
selv skulle undersøges. --- Saa at saadan Videnskab (Ontologie og Logik) altsaa ikke
kan træde op som mere forudsætningsløs, end Mathematiken. Derimod kan vistnok saadan
Videnskab, som blot explicativ, eller som gaaende ud fra simpel aandelig Iagttagelse
af de i al Tænkning fremtrædende Grundbestemmelser, træde frem med en vis
Forudsætningsløshed, ikke ulig den, man kan tillægge Mathematiken. Imidlertid vil dog
deri bestandigen et Hensyn til Virkeligheden blive at tage, som den rene Mathematik,
der kun bevæger sig i et reent formelt Muelighedernes Gebeet, ei har at tage. --- En
saadan blot explicativ Ontologie og Logik er det ogsaa, vi have for os i Hegels
saakaldte objective og subjective Logik, der imidlertid ikke allene vil træde op som
aldeles forudsætningsløs, men endog synes at skulle gjelde for det egentlige
Grundlag, hvilket den dog aldrig kan være, naar der skal menes det altomfattende
Grundlag, som kun Grundtrækkene til en altomfattende Verdensanskuelse kunne
give.\par}

% Rettelse (errata leaf, p.80 l.13): print reads ``Philosoohiens'' (dropped p); the
% erratum ``Philosophiens'' corrects it. Applied below (second occurrence, before
% ``Forhold til Aabenbaringen'').
{\small\emph{Anm.~7.} Paa dette Sted kunne vi nu og, med særdeles Hensyn til
Philosophiens Rettigheder, bemærke, hvad der er at kalde Speculationens Gebeet, at
dette nemlig er Eet med hele Livets og Tilværelsens Gebeet, naar der nemlig allene
sees hen til dettes rationale Grundlag; og at der altsaa, paa den ene Side, ikke kan
være noget særegent Indbegreb af Gjenstande, som kunde siges at høre Speculationen
allene til, men at der ogsaa, paa den anden Side, ei kan være Noget, som man vil
kunne udelukke Philosophien fra. Jeg bringer herved i Erindring, hvad jeg især med
Hensyn til Philosophiens Ret til og Opfordrethed til at gribe ind i den positiv
christelige Religion, ja i Christendommens hele Sag og dens hele Værk, har fremsat
deels i mit Philosoph.\ Archiv, saavel i første Hefte, især Pag.\ 71 ff., som i
fjerde Hefte, især Pag.\ 364 ff., deels i en tydsk Afhandling, som er trykt i tredie
Hefte af det af Schleiermacher, de Wette og Lücke i Aarene 1819--1822 udgivne
Theologische Zeitschrift, og angaaer Forholdet imellem christelig Tro og philosophisk
Erkjendelse. Paa disse Steder er iøvrigt Sagen ogsaa betragtet med Hensyn til
Christendommens og den Christnes egen Interesse, om jeg saa maa sige, saa at der ogsaa
er bleven handlet om Philosophiens Betydning for Christendommen, hvorved da tillige
Philosophiens Forhold til Aabenbaringen, som til en hellig Auctoritet, er bleven
omtalt. Iøvrigt vil jeg komme til atter at omhandle alt Dette i min
Christendomsphilosophie, naar jeg faaer den saavidt bragt til Fuldendelse, at jeg vil
kunne træde offentlig frem med den.\par}

\begin{center}\rule{0.18\linewidth}{0.4pt}\end{center}

%% ============================================================
%%  FULDELIG BESTEMMELSE AF PHILOSOPHIENS BEGREB.  (§ 20, p. 81; PDF 92)
%% ============================================================
\grouphead{Fuldelig Bestemmelse af Philosophiens Begreb.}

\begin{center}\large\S~20.\end{center}
\markboth{\S~20}{\S~20}
\begin{center}Endelig Definition paa Philosophie.\end{center}

Hovedmomenterne i Begrebet Philosophie kunne vi nu, forsaavidt Philosophien
betragtes reent objectivt, samle for os i følgende Definition:

\emph{Philosophie}, nemlig egentlig speculativ Philosophie, er den, under en
altomfattende Conferents af Alt med Alt og Discussion og Debat af Alt mod Alt, i
Kraft af Fornuft, altsaa udaf en apriorisk Grund --- den philosophiske Grundidee ---
sig udviklende Videnskab om alle Tings sande Væsen, ved en Tilbagegaaen til alle
Tings sidste --- eller, som de fra en anden Side vise sig at være, alle Tings første
--- Grunde og Grundforhold, og en Forfølgelse af alt Det, hvortil disses Erkjendelse
videre maa føre, saa at en altomfattende videnskabelig, altsaa systematisk,
Gudserkjendelse og Verdensanskuelse kan danne sig for os.

{\small\emph{Anm.} Ved Udtrykket: egentlig speculativ Philosophie tilkjendegives, at
Philosophie her tages \textit{sensu eminenti}, hvorved der da peges hen til, at ogsaa
dens Begreb har sit Dialektiske. --- Hvorvidt alle Philosopher vilde kunne bekjende
sig til denne vor endelige Definition, vil her ogsaa blive at overveie.\par}

\begin{center}\rule{0.18\linewidth}{0.4pt}\end{center}

%% ============================================================
%%  PHILOSOPHIEN SOM LED OG MOMENT I DEN HELE AANDELIGE LEVEN.
%%  (§ 21, pp. 82--86; PDF 93--97)
%% ============================================================
\grouphead{Philosophien som Led og Moment i den hele aandelige Leven.}

\begin{center}\large\S~21.\end{center}
\markboth{\S~21}{\S~21}
\begin{center}Philosophie paa lige Linie med det Empiriske og det Suprarationale,
--- med Poesie, Moral og Religion.\end{center}

At den hele Philosophie, om den end er en altomfattende Fornufterkjendelse, dog ikke
er en i alle Henseender altomfattende Erkjendelse, endsige at den skulde kunne være
en fuldelig Levens altomfattende Grundlag, maa allerede af det Foregaaende være
klart, da det deraf sees, at Det, vi philosophiskt erkjende, tillige maa sande sig
for os deels empiriskt deels sympathetiskt saa og gjennem den uafhængigt af
Speculationen, tilmed mere umiddelbart, sig uddannende aprioriske Rationalitet.
Allerede herved viser sig det Philosophiske kun at være et medbevægende Grundmoment
i Totaliteten, hvilket ei blot selv i sig har sin altomfattende Debat, men ogsaa med
alt sit hele Indhold i Henseende til Gyldigheden af sin hele Stræben, tilmed sin hele
Existents, drages ind i en større Debat, som ei længere føres paa Philosophiens Mark
allene, men paa den hele Videnskabeligheds Mark, hvor Philosophien da kun kan træde
frem som en Potents ved Siden af flere, ja endog maa søge at haandhæve sig i
Modsætning til, paa den ene Side Empirismen, paa den anden Side Suprarationalismen,
medens atter det Empiriske og det Suprarationale have at haandhæve sig i deres af al
Philosophie uafhængige Bestaaen.

Men endvidere fremkommer og subsisterer den hele Erkjenden selv, med den hele
Videnskabelighed overhovedet, atter kun i en Debat, som føres paa den hele Levens
Gebeet. Her bliver da den hele Videnskabelighed, med al den deri liggende Erkjenden,
nærmest og først en Potents lige overfor og i Samvirkning med Konsten, især den
ideelle Konst og den deri liggende Poesie, i hvilken den philosophiske
Verdensanskuelses ideelle Gehalt, altsaa Fornuftfylden, viser sig som udtrædende i og
virkeliggjørende sig i (incarnerende sig i) et ideelt Virkelighedens Gjenbillede.
Men dernæst staae begge disse Potentser endeligen i Forhold til Det i os, som gjør os
til Mere, end erkjendende og skuende Væsener, nemlig til følende og villende; og nu
maae baade Philosophie og Poesie, baade Videnskab og Konst, gaae op til Eenhed med
alt Dette, under en Debat, under hvilken det kommer an paa, at Gemyt og Villie fyldes
og styrkes ved dem, og atter virke tilbage paa dem; hvorunder da Religiøsitet, med
Tro og Kjærlighed, Tillid og Haab, i Gemytets, og Moralitet, med Ærefrygt og
Hellighed, i Villiens Gebeet, faae en Betydenhed lig den, som Philosophien har i
Videnskabelighedens og Poesien i Konstens Gebeet.

{\small\emph{Anm.~1.} Vi komme her atter tilbage til Modsætningen imellem Rationalt
og Extrarationalt og til den Bemærkning, at Liv og Tilvær bestaaer ved en væsentlig
Eenhed af begge (\S~13 b, Anm.~3, Pag.\ 50), hvortil vi kunne knytte den Bemærkning,
at Livet i Bevidstheden nærmest og først opstaaer i det Extrarationales Gebeet,
skjønt det ei kan røre sig i dette, uden strax at komme ind i Sammenhængs- og
Forholds-Erkjendelsens, altsaa i Rationalitetens Gebeet. Men nu bemærke vi, at til
den Rationalitet, som saaledes strax gjør sig gjeldende, komme vi nu enten fra det
Nærmestes Side, gjennem det ideelle Blik paa det umiddelbart Givne og Forfølgelsen af
hvad der nærmest knytter sig hertil, medens det Givne i dets Givethed bliver
Grundlaget, være det nu i Empiriens eller i den suprarationale Bevægetheds Gebeet;
eller ogsaa ved at gaae til de dybere Grunde og at gaae ud fra disse, for at udtyde
os dette Givne, saa at dets Forklaring og den saaledes uddannede Theorie bliver os
Grundlaget. Herpaa er det, at overhovedet Modsætningen imellem Rationalisme og
Empirisme i den ene, Modsætningen imellem Rationalisme og Suprarationalisme i den
anden Henseende beroer. Som Suprarationalt betragte vi her Alt, hvad der, ved en
anden Magt end Rationalitetens, samt ud over Det, vi see Vei til at gjøre os Rede
for, gjør sig saaledes gjeldende for os som noget Høiere, at vi ei kunne Andet, end
give det Rum. Dog bliver især kun Det, som i religiøs Henseende, eller med Hensyn til
den Livets Styrelse, vi maae give os hen i, udøver denne Magt over os, kaldet det
Suprarationale; hvilket da tillige, forsaavidt denne dets Magt gaaer ud fra et
Udspring, der kan ansees for at overstige den hele almindelige jordiske Tingenes
Orden (som punct- og glimtviis fremtrædende Frembrud af en høiere), kaldes det
Supranaturale. --- At vi bestandigen leve og røre os, som i det Empiriskes, saa og i
det Suprarationales Region (idet vi f.\ Ex.\ bestandigen leve i Tilliden til en
Styrelse), vil man let see.\par}

{\small\emph{Anm.~2.} Om Forholdet imellem Philosophie og Christendom, hvorom ovenfor
i Anm.~7.\ til \S~19, P.\ 79--80 er bleven talt, skal her ei videre handles, kun at
jeg her vil bringe i Erindring følgende Ord, som læses i mit Philos.\ Arch.\ Pag.\ 62:
„om end Philosophien er Aanden i den christelige Erkjenden, saa er dog den
tillidsfulde, ud over al Seen og Erkjenden gaaende, alt det Erkjendte og meget Mere,
end det, som i en Livskjærne, i sig indsluttende Tro, der, som en Hengivelse i
Overbeviisningens Magt i Liv og Gjerning, er en praktisk Leven i og med denne
Overbeviisnings Fylde, --- Sjælen i den hele christelige Existents, i hvilken hiint
Philosophiske kun er eet Moment --- men vistnok et væsentligt Grundmoment."\par}

{\small\emph{Anm.~3.} Under den philosophiske Stræbens Fremtræden i den større
Menneskeleven bliver det en naturlig Følge af det Sporadiske i Tilværelsen, at der
danne sig flere saakaldte Philosophier, der føre til Dannelsen af forskjellige
philosophiske Skoler, deels i theoretisk deels i praktisk Henseende.\par}

{\small\emph{Anm.~4.} Vi kunne her endnu tilsidst see hen til den Suprematie, som
Hegel og hans Skole har villet gjøre Fordring paa for Philosophien. For
Erkjendelsesbestræbelsen, i dens Fremgang til den fulde Sandheds grundige
Erkjendelse, staaer som høieste Formaal Erkjendelsen af Sandheden i dens Sandhed,
samt i Udfoldningen tillige af alle dens Momenter, som saadanne, og disses
Indoptagelse i Eenheden; og her træder da fuldendt Philosophie, som fuldelig
gjennemført Sandhedserkjendelse, frem som det Høieste, hvortil vi sigte.\par}

Men tænke vi os nu at være komne saavidt, maae vi dog see og erkjende, at
Sandhedsriget nu dog ogsaa maa vise sig for os i sin reelle Magt i selve det concrete
Liv, som det heri realiserede og bestandigen sig realiserende Ideerige. Men da see
vi, at vi, for at komme til det Fuldelige, ei kunne blive staaende ved Philosophien,
men maae ud over den. Hertil vil da nu først og fremmest henhøre Philosophiens
Opgaaen til Eenhed med vor Leven i Virkelighedens hele Sphære, forsaavidt vi paa
aldeles concret Maade ere og røre os heri, altsaa med vor Leven i Natur, Menneskeliv
og Historie. Men Dette gaaer nu videre; thi nu danner sig Konsten, nemlig den ideelle
Konst, som det Livets Gjenbillede, hvori vi see det Uendeliges Fylde i magtfuld
Nærværelse røre sig og leve i det concrete Livs hele Spil, hvorved det da tillige
tiltaler os saa, at det skues med Velbehag, saa at vi nu ganske anderledes
personligen indoptage det og leve os ind i det, end vi kunde, saalænge det kun stod
for os i sin Almeengyldighed. Og saaledes see vi da, at dersom Skjønhed var nok til
at betegne det hele æsthetiskt-poetiskt Tiltalende, kunde og maatte vi sige:
Sandhedens Rige skal være os Mere, end et Sandhedens Rige, det skal vorde, os et
Skjønhedens Rige, men bedre sige vi: et Poesiens Rige.

Men Mennesket skal ei blot studere det Algyldige som saadant, og fryde sig derved;
han maa og sætte sig i det rette Rapport eller Grundforhold dertil i Henseende til
sit personlige Selv, eller han maa gjøre det Algyldige til Algyldigt for ham selv, i
hans egen inderste Leven, saaledes som denne skal gaae frem af det personlige
Centrale i ham. Dette Centrale finde vi i Menneskets Villie, hvis væsentligste
Grundyttring maa være den personlige Attraaes og Tragtens Hengivelse i en algyldig
altomfattende Villie, der da constituerer sig for os, deels som Herskeren og Styreren
i al vor Tragten, Villen og Gjøren, deels som vort dybeste Tilhold og Gjenstand for
vor dybeste Tilknytning og Hengivelse, forsaavidt vi ere undergivne
Livstilskikkelsernes totale Gang. Saaledes ville vi da komme til at erkjende, at det
ikke er nok, at Sandhedens Rige gaaer op for os som det, der tillige træder frem som
et Skjønhedens, eller bedre et Poesiens Rige, men at det ogsaa maa vise sig for os
som det Godes Rige; men denne Betragtning fuldender sig først for os, naar vi nu,
idet vi gaae videre, see, at vi maae leve i Sandhedens og Livsfyldens Rige, som i
det, der tillige er os et Kjærlighedens og Hellighedens Rige, og vor herved betegnede
Indoptagethed i Livets Fylde vil da nu staae for os som det høieste Formaal.

\begin{center}\rule{0.28\linewidth}{0.4pt}\end{center}

%% ============================================================
%%  RETTELSER (errata leaf, PDF 11). Transcribed faithfully.
%%  NB items are substantive; folded into the body as it is transcribed.
%% ============================================================
%% Transcribed verbatim from the errata leaf (PDF 11). On the leaf the ``Pag.''
%% and ``Lin.'' values are given with ditto dashes for repeated page numbers; the
%% page/line numbers are here written out explicitly for clarity. ``l.'' = ``læs''.
\section*{Rettelser.}
\markboth{Rettelser}{Rettelser}

\noindent
\begin{tabular}{@{}l r r l@{}}
    & \multicolumn{1}{c}{Pag.} & \multicolumn{1}{c}{Lin.} & \\
    & 9  & 26     & bortfalder Ordet i.\\
    & 9  & 28     & for „ikke i selv" læs: „ikke i sig selv".\\
    & 14 & 28     & Tilfredsstillelelse, l.\ Tilfredsstillelse.\\
    & 18 & 3      & Tredelig, l.\ Tredeling.\\
    & 24 & 13     & snaledes, l.\ saaledes.\\
    & 25 & 11     & bemærkede, l.\ Bemærkede.\\
    & 26 & 16     & det, l.\ Dette.\\
    & 30 & 17     & bestemt, l.\ Bestemt.\\
    & 31 & 26     & opjective, l.\ objective.\\
NB. & 31 & 27     & altomfattende, l.\ altomfattende Maade.\\
NB. & 33 & 4      & Begrebets, l.\ Begrebet.\\
    & 35 & 20     & udslettes Kommaet.\\
NB. & 39 & 17--18 & for Theologie, l.\ Teleologie.\\
NB. & 40 & 18     & for det, l.\ de.\\
    & 41 & 14     & hvis, l.\ vis.\\
    & 43 & 24     & Jakobi, l.\ Jacobi.\\
    & 54 & 20     & saameget, l.\ saa meget.\\
NB. & 60 & 19     & sættes Komma efter Enkeltheder.\\
    & 61 & 15     & sættes Komma efter spørges.\\
    & 75 & 15     & læs: Tilværelsessphære.\\
    & 80 & 13     & Philosophiens.\\
\end{tabular}

\medskip
\noindent
Endnu bemærkes, at ved en Feiltagelse have to efter hinanden følgende Paragrapher
faaet Tallet 13; de ere i Indholdsangivelsen anførte som \S~13, a og \S~13, b.

\smallskip
\noindent
Ved denne Leilighed vil jeg gjøre opmærksom paa, at der i min Psychologie, ny
Udarbeidelse, Pag.\ 133, findes den slemme Feil, at der i nederste Linie staaer
umiddelbare, for: middelbare.

\end{document}
